% AdvancedMechanicsAssumptions20261023-v1
% AdvancedMechanicsQuality20261023-v1
\begingroup
\newif\ifAdvHideHeading
\RenewDocumentEnvironment{advexamplebox}{m}{%
  \refstepcounter{advexample}%
  \begin{tcolorbox}[
    enhanced,unbreakable,colback=white,colframe=RuleGray,boxrule=0.5pt,arc=0mm,
    title={例~\theadvexample\quad #1},fonttitle=\bfseries,
    coltitle=Ink,colbacktitle=PaleBlue,borderline west={1.4pt}{0pt}{DeepBlue}]
}{\end{tcolorbox}}

\RenewDocumentEnvironment{advsolution}{+b}{%
  \par\smallskip\begingroup\small\noindent
  \ifAdvHideHeading\global\AdvHideHeadingfalse\else\LectureSolutionHeading\fi #1\par\endgroup\smallskip
}{}

\RenewDocumentEnvironment{advdefinition}{m +b}{%
  \refstepcounter{advdefinition}%
  \begin{tcolorbox}[
    enhanced,unbreakable,colback=white,colframe=RuleGray,boxrule=0.5pt,arc=0mm,
    title={定义~\theadvdefinition\quad #1},fonttitle=\bfseries,
    coltitle=Ink,colbacktitle=PaleGold,borderline west={1.2pt}{0pt}{Gold}]
  #2
  \end{tcolorbox}
}{}

\RenewDocumentEnvironment{advtheorem}{m +b}{%
  \refstepcounter{advtheorem}%
  \begin{tcolorbox}[
    enhanced,unbreakable,colback=white,colframe=RuleGray,boxrule=0.5pt,arc=0mm,
    title={定理~\theadvtheorem\quad #1},fonttitle=\bfseries,
    coltitle=Ink,colbacktitle=PaleBlue,borderline west={1.2pt}{0pt}{DeepBlue}]
  #2
  \end{tcolorbox}
}{}


\AdvancedUnit{1}{力学}

\section{分析力学基础}

\subsection{广义坐标、约束与拉格朗日方程}

牛顿方程
\[
  m_i\ddot{\vect r}_i=\vect F_i,\qquad \dot{\vect p}_i=\vect F_i
\]
对直角坐标很直接，但在多体系统、曲线坐标和带约束系统中，它会引入大量未知约束力。
分析力学的做法是：先用适合体系几何结构的坐标消去约束，再用标量函数统一编码动力学。

设系统含有 $n$ 个质点，完整约束写成
\begin{equation}\label{advA1:eq:constraints}
 f^A(\vect r_1,\ldots,\vect r_n,t)=0,
 \qquad A=1,\ldots,N .
\end{equation}
若它们相互独立，系统自由度为 $f=3n-N$。对应的 $f$维可运动空间称为位形空间。
广义坐标 $q_1,\ldots,q_f$ 的选取应使映射
\begin{equation}\label{advA1:eq:point-map}
 \vect r_i=\vect r_i(q_1,\ldots,q_f,t)
\end{equation}
自动满足式\eqref{advA1:eq:constraints}。约束于是不再以附加方程出现，而是被吸收到函数\eqref{advA1:eq:point-map}的结构中。不显含时间的约束称为稳定约束；显含时间的称为非稳定约束。本节不处理无法积分为位形约束的非完整约束。

\begin{advexamplebox}{刚性椭球的十二个约束}\label{advA1:ex:ellipsoid}
六个质点固定在椭球的六个顶点 $A,A',B,B',C,C'$ 上，
$AA'=2a,BB'=2b,CC'=2c$。用十八个直角坐标写出一组相互独立的十二个约束。

\begin{qualityfigure}
\begin{tikzpicture}[>=Latex,font=\small]
\draw[DeepBlue](0,0)ellipse(2.5 and 1.5);\draw[dashed](-2.5,0)arc(180:360:2.5 and .45);\draw(-2.5,0)arc(180:0:2.5 and .45);
\foreach\x/\y/\t in {2.5/0/A,-2.5/0/A',.9/.45/B,-.9/-.45/B',0/1.5/C,0/-1.5/C'}{\fill(\x,\y)circle(2pt)node[above right]{$\t$};}
\draw[dashed](-2.5,0)--(2.5,0)(-.9,-.45)--(.9,.45)(0,-1.5)--(0,1.5);\draw[<->](-2.5,-2.1)--(2.5,-2.1)node[midway,below]{$AA'=2a$};\draw[<->](-3,-1.5)--(-3,1.5)node[midway,left]{$CC'=2c$};\node at(1.7,.85){$BB'=2b$};
\end{tikzpicture}
\figcap{六个对径顶点和三条主轴（题设几何）}\label{advA1:fig:quality-ellipsoid-q}
\end{qualityfigure}
\end{advexamplebox}
\begin{advsolution}
\begin{qualityfigure}
\begin{tikzpicture}[>=Latex,font=\small]
\draw[DeepBlue](0,0)ellipse(2.5 and 1.5);\draw[dashed](-2.5,0)arc(180:360:2.5 and .45);\draw(-2.5,0)arc(180:0:2.5 and .45);
\foreach\x/\y/\t in {2.5/0/A,-2.5/0/A',.9/.45/B,-.9/-.45/B',0/1.5/C,0/-1.5/C'}{\fill(\x,\y)circle(2pt)node[above right]{$\t$};}
\draw[dashed](-2.5,0)--(2.5,0)(-.9,-.45)--(.9,.45)(0,-1.5)--(0,1.5);\draw[<->](-2.5,-2.1)--(2.5,-2.1)node[midway,below]{$AA'=2a$};\draw[<->](-3,-1.5)--(-3,1.5)node[midway,left]{$CC'=2c$};\node at(1.7,.85){$BB'=2b$};\node[align=left]at(6.5,0){三组对径端点共用同一中点\\$|\vect u|=2a,\ |\vect v|=2b,\ |\vect w|=2c$\\三条对径矢量两两正交};
\end{tikzpicture}
\figcap{六个对径顶点和三条主轴（解答辅助图）}\label{advA1:fig:quality-ellipsoid-a}
\end{qualityfigure}
图~\ref{advA1:fig:quality-ellipsoid-a}中，$A,A'$、$B,B'$、$C,C'$ 各自的中点必须相同；长度条件与正交条件是另外六个标量约束，不能把六个端点各自固定在实验室坐标中。

引入三条对径矢量
\[
 \vect u=\vect r_A-\vect r_{A'},\quad
 \vect v=\vect r_B-\vect r_{B'},\quad
 \vect w=\vect r_C-\vect r_{C'}.
\]
先要求三条对径共中点：
\[
 \vect r_A+\vect r_{A'}-\vect r_B-\vect r_{B'}=\vect0,
 \quad
 \vect r_A+\vect r_{A'}-\vect r_C-\vect r_{C'}=\vect0 .
\]
这是 $6$ 个标量方程。再规定对径长度和两两正交：
\[
 \vect u^2=4a^2,\qquad\vect v^2=4b^2,\qquad\vect w^2=4c^2,
 \qquad
 \vect u\!\cdot\!\vect v=\vect v\!\cdot\!\vect w
 =\vect w\!\cdot\!\vect u=0 .
\]
共有 $6+3+3=12$ 个独立标量约束，因而自由度为 $18-12=6$，与一般刚体的三个平动和三个转动自由度一致。
\end{advsolution}


位形空间中的光滑轨道 $q_\alpha=q_\alpha(t)$ 确定了系统的运动。一般动力学量是
$A(q_\alpha,\dot q_\alpha,t)$，沿轨道的全导数为
\begin{equation}\label{advA1:eq:total-derivative}
 \dot A=\sum_\alpha\frac{\partial A}{\partial q_\alpha}\dot q_\alpha
 +\sum_\alpha\frac{\partial A}{\partial\dot q_\alpha}\ddot q_\alpha
 +\frac{\partial A}{\partial t}.
\end{equation}
由式\eqref{advA1:eq:point-map}，
\[
 \dot{\vect r}_i=\sum_\alpha
 \frac{\partial\vect r_i}{\partial q_\alpha}\dot q_\alpha
 +\frac{\partial\vect r_i}{\partial t}.
\]
因此动能对广义速度总是二次多项式：
\begin{align}
 T&=T_2+T_1+T_0,\label{advA1:eq:T-decomp}\\
 T_2&=\frac12\sum_{\alpha\beta}m_{\alpha\beta}\dot q_\alpha\dot q_\beta,
 &m_{\alpha\beta}&=\sum_i m_i
 \frac{\partial\vect r_i}{\partial q_\alpha}\!\cdot\!
 \frac{\partial\vect r_i}{\partial q_\beta},\notag\\
 T_1&=\sum_\alpha\pi_\alpha\dot q_\alpha,
 &\pi_\alpha&=\sum_i m_i
 \frac{\partial\vect r_i}{\partial q_\alpha}\!\cdot\!
 \frac{\partial\vect r_i}{\partial t},\notag\\
 T_0&=\frac12\tau,
 &\tau&=\sum_i m_i\left(\frac{\partial\vect r_i}{\partial t}\right)^2.\notag
\end{align}
稳定约束且坐标变换不含时时，$T_1=T_0=0$。

\begin{advdefinition}{广义机械动量}\label{advA1:def:momentum}
\[
 p_\alpha^{(\mathrm{mech})}=\frac{\partial T}{\partial\dot q_\alpha}
 =\sum_\beta m_{\alpha\beta}\dot q_\beta+\pi_\alpha .
\]
\end{advdefinition}

\begin{advdefinition}{广义力}\label{advA1:def:force}
对虚位移 $\dd\vect r_i=\sum_\alpha
(\partial\vect r_i/\partial q_\alpha)\dd q_\alpha$，定义
\[
 Q_\alpha=\sum_i\vect F_i\!\cdot\!
 \frac{\partial\vect r_i}{\partial q_\alpha}.
\]
若坐标变换显含时间，实际元功还包含
\[
 \dd W=\sum_\alpha Q_\alpha\dd q_\alpha+h\dd t,
 \qquad h=\sum_i\vect F_i\!\cdot\!\frac{\partial\vect r_i}{\partial t}.
\]
\end{advdefinition}

\begin{advexamplebox}{在膨胀球面上运动的质点}\label{advA1:ex:expanding-sphere}
设 $v>0,t>0$。质量为 $m$ 的质点被约束在 $x^2+y^2+z^2-v^2t^2=0$ 上。以球坐标角
$\theta,\varphi$ 为广义坐标，求位矢、动能、广义动量及其全导数、广义力和元功。

\begin{qualityfigure}
\begin{tikzpicture}[>=Latex,font=\small]
\draw[->](0,0)--(3.5,0)node[right]{$x$};\draw[->](0,0)--(0,3)node[above]{$z$};\draw[->](0,0)--(-1.3,-1)node[left]{$y$};\node[below left]at(0,0){$O$};\draw[DeepBlue](0,0)circle(2);\coordinate(P)at(1.3,1.3);\fill(P)circle(2pt)node[right]{$P$};\draw[thick](0,0)--(P)node[pos=.65,below right=3pt]{$r$};
\draw[dashed](P)--(1.3,-.45)--(0,0);\draw(0,.7)arc(90:45:.7)node[midway,above]{$\theta$};\draw(.7,0)arc(0:-19.1:.7)node[midway,below]{$\varphi$};\node at(3.5,2){$r(t)=vt$};
\end{tikzpicture}
\figcap{膨胀球面的瞬时半径和两角坐标（题设几何）}\label{advA1:fig:quality-expanding-sphere-q}
\end{qualityfigure}
\end{advexamplebox}
\begin{advsolution}
\begin{qualityfigure}
\begin{tikzpicture}[>=Latex,font=\small]
\draw[->](0,0)--(3.5,0)node[right]{$x$};\draw[->](0,0)--(0,3)node[above]{$z$};\draw[->](0,0)--(-1.3,-1)node[left]{$y$};\node[below left]at(0,0){$O$};\draw[DeepBlue](0,0)circle(2);\coordinate(P)at(1.3,1.3);\fill(P)circle(2pt)node[right]{$P$};\draw[thick](0,0)--(P)node[pos=.65,below right=3pt]{$r$};
\draw[dashed](P)--(1.3,-.45)--(0,0);\draw(0,.7)arc(90:45:.7)node[midway,above]{$\theta$};\draw(.7,0)arc(0:-19.1:.7)node[midway,below]{$\varphi$};\node at(3.5,2){$r(t)=vt$};\draw[->,Gold](P)--(2.1,2.1)node[right]{$\vect e_r$};\draw[->,Gold](P)--(2.1,.5)node[right]{$\vect e_\theta$};\node[align=left]at(4,-1.4){$\dd s_r=\dd r$\\$\dd s_\theta=r\dd\theta$\\$\dd s_\varphi=r\sin\theta\dd\varphi$};
\end{tikzpicture}
\figcap{膨胀球面的瞬时半径和两角坐标（解答辅助图）}\label{advA1:fig:quality-expanding-sphere-a}
\end{qualityfigure}
图~\ref{advA1:fig:quality-expanding-sphere-a}中，$P$ 的径向距离已由非稳定约束固定为 $vt$；仍可独立变化的是 $\theta,\varphi$，而实际微位移还含球面膨胀的径向部分。

\[
 \vect r=vt(\sin\theta\cos\varphi,\sin\theta\sin\varphi,\cos\theta),
\]
\[
 T=\frac m2\left[v^2+v^2t^2
 (\dot\theta^2+\sin^2\theta\,\dot\varphi^2)\right].
\]
于是
\[
 p_\theta=mv^2t^2\dot\theta,
 \qquad p_\varphi=mv^2t^2\sin^2\theta\,\dot\varphi,
\]
\[
 \dot p_\theta=mv^2(2t\dot\theta+t^2\ddot\theta),
\]
\[
 \dot p_\varphi=mv^2\!\left[
 2t\sin^2\theta\dot\varphi
 +2t^2\sin\theta\cos\theta\dot\theta\dot\varphi
 +t^2\sin^2\theta\ddot\varphi\right].
\]
若合力的球坐标分量为 $(F_r,F_\theta,F_\varphi)$，则
\[
 Q_\theta=vtF_\theta,\qquad Q_\varphi=vt\sin\theta F_\varphi,
 \qquad h=vF_r,
\]
\[
 \dd W=Q_\theta\dd\theta+Q_\varphi\dd\varphi+vF_r\dd t.
\]
最后一项是球面本身膨胀方向上的做功，不能误并入两个角向广义力。
\end{advsolution}

\begin{advtheorem}{拉格朗日方程的引理}\label{advA1:thm:lagrange-lemma}
沿任一满足牛顿方程的运动，
\begin{equation}\label{advA1:eq:lagrange-lemma}
 \frac{\dd}{\dd t}\frac{\partial T}{\partial\dot q_\alpha}
 -\frac{\partial T}{\partial q_\alpha}=Q_\alpha .
\end{equation}
\end{advtheorem}
证明不再省略。由链式法则，
\[
 \frac{\partial T}{\partial\dot q_\alpha}
 =\sum_i m_i\dot{\vect r}_i\!\cdot\!
 \frac{\partial\vect r_i}{\partial q_\alpha}.
\]
对时间求全导数，并注意
$\partial\dot{\vect r}_i/\partial q_\alpha
=\dd(\partial\vect r_i/\partial q_\alpha)/\dd t$，得
\[
 \frac{\dd}{\dd t}\frac{\partial T}{\partial\dot q_\alpha}
 =\sum_i m_i\ddot{\vect r}_i\!\cdot\!
 \frac{\partial\vect r_i}{\partial q_\alpha}
 +\sum_i m_i\dot{\vect r}_i\!\cdot\!
 \frac{\partial\dot{\vect r}_i}{\partial q_\alpha}.
\]
第二项正是 $\partial T/\partial q_\alpha$，两式相减，再用
$m_i\ddot{\vect r}_i=\vect F_i$ 便得式\eqref{advA1:eq:lagrange-lemma}。

\begin{advexamplebox}{二维极坐标中的动力学方程}\label{advA1:ex:polar}
由拉格朗日引理推导平面极坐标中的质点动力学方程。

\begin{qualityfigure}
\begin{tikzpicture}[>=Latex,font=\small]
\draw[->](0,0)--(4.6,0)node[right]{$x$};\draw[->](0,0)--(0,3.4)node[above]{$y$};\node[below left]at(0,0){$O$};\coordinate(P)at(2.5,1.8);\fill(P)circle(2pt)node[left]{$P$};\draw[thick](0,0)--(P)node[midway,below]{$r$};\draw(.9,0)arc(0:35.75:.9)node[midway,right]{$\theta$};
\end{tikzpicture}
\figcap{极坐标中随位置旋转的正交基矢（题设几何）}\label{advA1:fig:quality-polar-q}
\end{qualityfigure}
\end{advexamplebox}
\begin{advsolution}
\begin{qualityfigure}
\begin{tikzpicture}[>=Latex,font=\small]
\draw[->](0,0)--(4.6,0)node[right]{$x$};\draw[->](0,0)--(0,3.4)node[above]{$y$};\node[below left]at(0,0){$O$};\coordinate(P)at(2.5,1.8);\fill(P)circle(2pt)node[left]{$P$};\draw[thick](0,0)--(P)node[midway,below]{$r$};\draw(.9,0)arc(0:35.75:.9)node[midway,right]{$\theta$};\draw[->,DeepBlue](P)--(3.5,2.52)node[right]{$\vect e_r$};\draw[->,DeepBlue](P)--(1.78,2.8)node[left]{$\vect e_\theta$};
\end{tikzpicture}
\figcap{极坐标中随位置旋转的正交基矢（解答辅助图）}\label{advA1:fig:quality-polar-a}
\end{qualityfigure}
图~\ref{advA1:fig:quality-polar-a}中，径向与横向单位矢量随 $\theta$ 改变。先把速度写成径向、横向两项，再求导时须同时对基矢求导。

$T=m(\dot r^2+r^2\dot\theta^2)/2$，而 $Q_r=F_r, Q_\theta=rF_\theta$。代入式\eqref{advA1:eq:lagrange-lemma}：
\[
 m(\ddot r-r\dot\theta^2)=F_r,
 \qquad
 m(r\ddot\theta+2\dot r\dot\theta)=F_\theta.
\]
第二式括号内每项均为加速度；乘以质量 $m$ 后与右端力同量纲。若漏掉 $r\ddot\theta$ 中的 $r$，量纲会立即失配。
\end{advsolution}

\begin{advexamplebox}{三维球坐标中的动力学方程}\label{advA1:ex:spherical}
推导球坐标 $(r,\theta,\varphi)$ 中的三个动力学方程，并与加速度的矢量分解比较。

\begin{qualityfigure}
\begin{tikzpicture}[>=Latex,font=\small]
\draw[->](0,0)--(3.5,0)node[right]{$x$};\draw[->](0,0)--(0,3)node[above]{$z$};\draw[->](0,0)--(-1.3,-1)node[left]{$y$};\node[below left]at(0,0){$O$};\draw[DeepBlue](0,0)circle(2);\coordinate(P)at(1.3,1.3);\fill(P)circle(2pt)node[right]{$P$};\draw[thick](0,0)--(P)node[pos=.65,below right=3pt]{$r$};
\draw[dashed](P)--(1.3,-.45)--(0,0);\draw(0,.7)arc(90:45:.7)node[midway,above]{$\theta$};\draw(.7,0)arc(0:-19.1:.7)node[midway,below]{$\varphi$};
\end{tikzpicture}
\figcap{球坐标的径向、极角和方位角位移（题设几何）}\label{advA1:fig:quality-spherical-q}
\end{qualityfigure}
\end{advexamplebox}
\begin{advsolution}
\begin{qualityfigure}
\begin{tikzpicture}[>=Latex,font=\small]
\draw[->](0,0)--(3.5,0)node[right]{$x$};\draw[->](0,0)--(0,3)node[above]{$z$};\draw[->](0,0)--(-1.3,-1)node[left]{$y$};\node[below left]at(0,0){$O$};\draw[DeepBlue](0,0)circle(2);\coordinate(P)at(1.3,1.3);\fill(P)circle(2pt)node[right]{$P$};\draw[thick](0,0)--(P)node[pos=.65,below right=3pt]{$r$};
\draw[dashed](P)--(1.3,-.45)--(0,0);\draw(0,.7)arc(90:45:.7)node[midway,above]{$\theta$};\draw(.7,0)arc(0:-19.1:.7)node[midway,below]{$\varphi$};\draw[->,Gold](P)--(2.1,2.1)node[right]{$\vect e_r$};\draw[->,Gold](P)--(2.1,.5)node[right]{$\vect e_\theta$};\node[align=left]at(4,-1.4){$\dd s_r=\dd r$\\$\dd s_\theta=r\dd\theta$\\$\dd s_\varphi=r\sin\theta\dd\varphi$};
\end{tikzpicture}
\figcap{球坐标的径向、极角和方位角位移（解答辅助图）}\label{advA1:fig:quality-spherical-a}
\end{qualityfigure}
图~\ref{advA1:fig:quality-spherical-a}中，三个线元分别为 $\dd r,r\dd\theta,r\sin\theta\dd\varphi$；第三项的纬圈半径是 $r\sin\theta$，不能漏掉这个度量因子。

\[
 T=\frac m2[\dot r^2+r^2\dot\theta^2+r^2\sin^2\theta\dot\varphi^2],
\quad
 (Q_r,Q_\theta,Q_\varphi)=(F_r,rF_\theta,r\sin\theta F_\varphi).
\]
逐个代入引理得
\begin{align*}
 m[\ddot r-r(\dot\theta^2+\sin^2\theta\dot\varphi^2)]&=F_r,\\
 m[r\ddot\theta+2\dot r\dot\theta-r\sin\theta\cos\theta\dot\varphi^2]&=F_\theta,\\
 m[r\sin\theta\ddot\varphi+2\dot r\sin\theta\dot\varphi
 +2r\cos\theta\dot\theta\dot\varphi]&=F_\varphi.
\end{align*}
直接对 $\vect r=r\vect e_r$ 求两次导数，再投影到
$(\vect e_r,\vect e_\theta,\vect e_\varphi)$，会得到完全相同的三个分量，这正是坐标协变性的具体检验。
\end{advsolution}

\begin{advexamplebox}{平移且转动的参考系}\label{advA1:ex:moving-frame}
建立坐标变换
\[
 \binom{\bar x}{\bar y}=
 \begin{pmatrix}\cos\theta&\sin\theta\\-\sin\theta&\cos\theta\end{pmatrix}
 \left[\binom{x}{y}-\binom{X}{Y}\right],
\]
其中 $X,Y,\theta$ 均为时间函数。解释其几何意义，写出元功并推导新系中的运动方程。

\begin{qualityfigure}
\begin{tikzpicture}[>=Latex,font=\small]
\draw[->](0,0)--(4.2,0)node[right]{$x$};\draw[->](0,0)--(0,3.2)node[above]{$y$};\node[below left]at(0,0){$O$};\coordinate(C)at(1.8,1);\fill(C)circle(2pt)node[below right]{$O'$};\draw[->,Gold](0,0)--(C)node[midway,below]{$\vect R$};\draw[->,DeepBlue](C)--(3.8,2.2)node[right]{$\bar x$};\draw[->,DeepBlue](C)--(.8,2.67)node[above]{$\bar y$};\draw(2.4,1)arc(0:31:.6)node[midway,right]{$\theta$};
\end{tikzpicture}
\figcap{原点平移与坐标基矢转动（题设几何）}\label{advA1:fig:quality-moving-frame-q}
\end{qualityfigure}
\end{advexamplebox}
\begin{advsolution}
\begin{qualityfigure}
\begin{tikzpicture}[>=Latex,font=\small]
\draw[->](0,0)--(4.2,0)node[right]{$x$};\draw[->](0,0)--(0,3.2)node[above]{$y$};\node[below left]at(0,0){$O$};\coordinate(C)at(1.8,1);\fill(C)circle(2pt)node[below right]{$O'$};\draw[->,Gold](0,0)--(C)node[midway,below]{$\vect R$};\draw[->,DeepBlue](C)--(3.8,2.2)node[right]{$\bar x$};\draw[->,DeepBlue](C)--(.8,2.67)node[above]{$\bar y$};\draw(2.4,1)arc(0:31:.6)node[midway,right]{$\theta$};\node[align=left]at(4,-1.2){平移原点后再旋转分量\\相同物理矢量使用不同基矢表示};
\end{tikzpicture}
\figcap{原点平移与坐标基矢转动（解答辅助图）}\label{advA1:fig:quality-moving-frame-a}
\end{qualityfigure}
图~\ref{advA1:fig:quality-moving-frame-a}中，先减去原点位移 $\vect R=(X,Y)$，再将同一个相对位置投影到转过 $\theta$ 的新轴上。对坐标变换求导时，平移和基矢转动均随时间变化。

这是原点位于 $\vect R=(X,Y)$、坐标轴相对惯性系转过 $\theta$ 的非惯性系。令
$\vect\rho=(\bar x,\bar y)$，$\vect\Omega=\dot\theta\vect e_z$，则
\[
 \vect r=\vect R+\mathsf R(\theta)\vect\rho,
\]
\[
 \dd W=\vect F'\!\cdot\!\dd\vect\rho
 +\vect F\!\cdot\![\dot{\vect R}+\vect\Omega\times
 \mathsf R(\theta)\vect\rho]\dd t,
\]
其中 $\vect F'=\mathsf R^{\mathsf T}\vect F$。两次求导得
\[
 m\ddot{\vect\rho}=\vect F'-m\vect a_0'
 -2m\vect\Omega\times\dot{\vect\rho}
 -m\dot{\vect\Omega}\times\vect\rho
 -m\vect\Omega\times(\vect\Omega\times\vect\rho).
\]
右边依次是真实力、平移惯性力、科里奥利力、欧拉力和离心力。所有项都是力，量纲为 $MLT^{-2}$。
\end{advsolution}

\begin{advdefinition}{理想约束}\label{advA1:def:ideal-constraint}
若约束 $f^A=0$ 对第 $i$ 个质点产生的约束力为
\[
 \vect P_i^{,A}=\lambda_A\nabla_i f^A,
\]
则称其为理想约束。由于
$\sum_i\nabla_i f^A\cdot\partial\vect r_i/\partial q_\alpha=0$，它对任一允许的虚位移都不做虚功，所以其广义力恒为零。
\end{advdefinition}

\begin{advdefinition}{广义势}\label{advA1:def:generalized-potential}
若存在 $V(q_\alpha,\dot q_\alpha,t)$ 使主动力的广义力满足
\[
 Q_\alpha=\frac{\dd}{\dd t}\frac{\partial V}{\partial\dot q_\alpha}
 -\frac{\partial V}{\partial q_\alpha},
\]
则称相互作用可由广义势描述。若 $V$ 不含速度，就回到 $Q_\alpha=-\partial V/\partial q_\alpha$。
\end{advdefinition}

\begin{advdefinition}{拉格朗日量}\label{advA1:def:lagrangian}
\[
 L(q_\alpha,\dot q_\alpha,t)=T(q_\alpha,\dot q_\alpha,t)
 -V(q_\alpha,\dot q_\alpha,t),
 \qquad p_\alpha=\frac{\partial L}{\partial\dot q_\alpha}.
\]
含速势时，$p_\alpha$ 是正则动量，一般不等于机械动量。
\end{advdefinition}

\begin{advtheorem}{拉格朗日方程}\label{advA1:thm:lagrange}
对可由广义势描述的理想约束系统，
\begin{equation}\label{advA1:eq:lagrange}
 \frac{\dd}{\dd t}\frac{\partial L}{\partial\dot q_\alpha}
 -\frac{\partial L}{\partial q_\alpha}=0,
 \qquad \alpha=1,\ldots,f .
\end{equation}
\end{advtheorem}
证明只需把定义\ref{advA1:def:generalized-potential}代入引理\ref{advA1:thm:lagrange-lemma}，再将动能项和势函数项合并。这一步也说明，方程的普遍性不是额外假设，而是由牛顿方程和点变换的链式法则得到的。

\begin{advexamplebox}{收短摆绳的圆锥摆}\label{advA1:ex:conical-pendulum}
均匀重力场中，质量 $m$ 的摆球以竖直向下方向为极轴，摆角为 $\theta$、方位角为 $\varphi$。摆长以常速 $v$ 收短，$l(t)=l_0-vt$。写出 $\theta,\varphi$ 的运动方程。

\begin{qualityfigure}
\begin{tikzpicture}[>=Latex,font=\small]
\fill(0,2.5)circle(2pt)node[above]{$O$};\draw[dashed](0,2.5)--(0,-1);\coordinate(P)at(1.5,.1);\draw[thick](0,2.5)--(P);\fill(P)circle(3pt)node[right]{$m$};\draw(0,1.6)arc(-90:-58:.9)node[midway,below right]{$\theta$};\draw[dashed](0,.1)ellipse(1.5 and .4);\node[above right]at(.75,1.3){$l(t)=l_0-vt$};
\end{tikzpicture}
\figcap{收短摆绳的摆角和瞬时轨道（题设几何）}\label{advA1:fig:quality-conical-pendulum-q}
\end{qualityfigure}
\end{advexamplebox}
\AdvHideHeadingtrue
\begin{advsolution}题图~\ref{advA1:fig:quality-conical-pendulum-q}中，摆长沿绳方向改变，方位角对应的水平轨道半径为 $l\sin\theta$；因而 $\varphi$ 的广义动量含 $l^2\sin^2\theta$，收绳时不能视为常系数。

\noindent\begin{minipage}{\linewidth}{\bfseries\color{DeepBlue}解答}\par
\begin{center}\begin{minipage}{\linewidth}\centering
\begin{tikzpicture}[scale=.95]
 \coordinate (O) at (0,1.5);\coordinate (P) at (1.3,-.5);
 \draw[axis] (0,1.8)--(0,-1.5) node[below]{$z$};
 \draw[guide] (0,-.5) ellipse (1.3 and .38);
 \draw[curve] (O)--(P) node[midway,right=5pt]{$l(t)$};
 \fill (O) circle(1.5pt) node[left]{$O$};\fill (P) circle(2.2pt) node[right=5pt]{$m$};
 \draw[DeepBlue,->] (0,.85) arc[start angle=-90,end angle=-56.98,radius=.65];
 \node at (.32,.70) {$\theta$};
 \draw[TextGray,->] ({1.3*cos(220)},{-.5+.38*sin(220)}) arc[start angle=220,end angle=300,x radius=1.3,y radius=.38];
 \node at (-.85,-1.08) {$\varphi$ 正向};
\end{tikzpicture}
\bookcaption{收短摆绳时的两个广义坐标}\label{advA1:fig:conical}
\end{minipage}\end{center}\end{minipage}

如图\ref{advA1:fig:conical}，
\[
 L=\frac m2[\dot l^2+l^2\dot\theta^2+l^2\sin^2\theta\dot\varphi^2]
 +mgl\cos\theta .
\]
$l$ 是给定函数而非待求自由度。对 $\theta,\varphi$ 列方程：
\[
 l^2\ddot\theta+2l\dot l\dot\theta
 -l^2\sin\theta\cos\theta\dot\varphi^2+gl\sin\theta=0,
\]
\[
 \frac{\dd}{\dd t}(ml^2\sin^2\theta\dot\varphi)=0,
 \qquad \dot l=-v.
\]
第二式表明绕竖直轴的角动量守恒；收绳装置可以交换能量，所以摆球的机械能一般不守恒。
\end{advsolution}

\begin{advexamplebox}{缓慢变化引力常量的开普勒问题}\label{advA1:ex:kepler-time}
平面有心力 $F_r=-mk(t)/r^2$。忽略使 $k$ 改变的物质流、辐射等其他作用，写出运动方程并判断哪些守恒量仍然存在。

\begin{qualityfigure}
\begin{tikzpicture}[>=Latex,font=\small]
\draw[->](0,0)--(4.6,0)node[right]{$x$};\draw[->](0,0)--(0,3.4)node[above]{$y$};\node[below left]at(0,0){$O$};\coordinate(P)at(2.5,1.8);\fill(P)circle(2pt)node[left]{$P$};\draw[thick](0,0)--(P)node[midway,below]{$r$};\draw(.9,0)arc(0:35.75:.9)node[midway,right]{$\theta$};\draw[->,Gold,very thick](P)--(1.7,1.224)node[above left]{$F_r$};\node at(3,-1){$k=k(t)$};
\end{tikzpicture}
\figcap{随时间变化的中心引力和极径（题设几何）}\label{advA1:fig:quality-kepler-time-q}
\end{qualityfigure}
\end{advexamplebox}
\begin{advsolution}
\begin{qualityfigure}
\begin{tikzpicture}[>=Latex,font=\small]
\draw[->](0,0)--(4.6,0)node[right]{$x$};\draw[->](0,0)--(0,3.4)node[above]{$y$};\node[below left]at(0,0){$O$};\coordinate(P)at(2.5,1.8);\fill(P)circle(2pt)node[left]{$P$};\draw[thick](0,0)--(P)node[midway,below]{$r$};\draw(.9,0)arc(0:35.75:.9)node[midway,right]{$\theta$};\draw[->,Gold,very thick](P)--(1.7,1.224)node[above left]{$F_r$};\node at(3,-1){$k=k(t)$};\draw[->,DeepBlue](P)--(3.5,2.52)node[right]{$\vect e_r$};\draw[->,DeepBlue](P)--(1.78,2.8)node[left]{$\vect e_\theta$};
\end{tikzpicture}
\figcap{随时间变化的中心引力和极径（解答辅助图）}\label{advA1:fig:quality-kepler-time-a}
\end{qualityfigure}
图~\ref{advA1:fig:quality-kepler-time-a}中，力始终沿 $PO$，所以对 $O$ 的力矩为零；但力的强度显含时间，角动量守恒与能量是否守恒须分别判断。

\[
 L=\frac m2(\dot r^2+r^2\dot\varphi^2)+\frac{mk(t)}r .
\]
拉格朗日方程给出
\[
 \ddot r-r\dot\varphi^2=-\frac{k(t)}{r^2},
 \qquad \frac{\dd}{\dd t}(mr^2\dot\varphi)=0.
\]
因为 $\varphi$ 仍是循环坐标，角动量严格守恒。哈密顿量
\[
 H=\frac m2(\dot r^2+r^2\dot\varphi^2)-\frac{mk(t)}r
\]
却满足
\[
 \dot H=-\frac{\partial L}{\partial t}=-\frac{m\dot k(t)}r,
\]
所以 $\dot k\ne0$ 时能量不守恒。量纲检查：$k$ 的量纲是 $L^3T^{-2}$，$mk/r$ 的量纲正是能量。
\end{advsolution}

\Needspace{12\baselineskip}
\subsection{守恒量、对称性与罗斯约化}

\begin{advdefinition}{哈密顿量}\label{advA1:def:hamiltonian}
\begin{equation}\label{advA1:eq:H}
 H(q,\dot q,t)=\sum_\alpha p_\alpha\dot q_\alpha-L,
 \qquad p_\alpha=\frac{\partial L}{\partial\dot q_\alpha}.
\end{equation}
在稳定约束、不含时坐标变换且 $V=V(q)$ 时，$H=T+V$。一般情形下不应把哈密顿量无条件地称为机械能。
\end{advdefinition}

\begin{advtheorem}{能量积分}\label{advA1:thm:energy}
沿真实运动，
\begin{equation}\label{advA1:eq:Hdot}
 \dot H=-\frac{\partial L}{\partial t}.
\end{equation}
因而当且仅当轨道上 $\partial L/\partial t=0$ 时，$H$ 守恒。
\end{advtheorem}
证明：对式\eqref{advA1:eq:H}求导，
\[
 \dot H=\sum_\alpha(\dot p_\alpha\dot q_\alpha+p_\alpha\ddot q_\alpha)
 -\sum_\alpha\left(\frac{\partial L}{\partial q_\alpha}\dot q_\alpha
 +\frac{\partial L}{\partial\dot q_\alpha}\ddot q_\alpha\right)
 -\frac{\partial L}{\partial t}.
\]
用 $\dot p_\alpha=\partial L/\partial q_\alpha$ 和
$p_\alpha=\partial L/\partial\dot q_\alpha$，前两组项逐项消去，即得式\eqref{advA1:eq:Hdot}。

\begin{advexamplebox}{系数随时间同比增长的谐振子}\label{advA1:ex:scaled-oscillator}
设 $m,k>0$；涉及缓变极限时取 $0<\gamma<\sqrt{k/m}$，并固定不随 $\gamma$ 改变的有限初始位置与速度。对 $L=(m\dot q^2-kq^2)/2$：
(1) 作坐标标度变换使动能系数为 $1$；
(2) 再考虑
\[
 L=\frac12me^{2\gamma t}\dot q^2-\frac12ke^{2\gamma t}q^2,
 \qquad \omega^2=k/m,
\]
求方程、哈密顿量及其全导数；
(3) 用 $\bar q=e^{\gamma t}q$ 找出守恒量；
(4) 证明在 $t=1/\gamma$ 的时间尺度上，$\gamma\to0$ 时哈密顿量的变化趋于零。

\begin{qualityfigure}
\begin{tikzpicture}[>=Latex,font=\small]
\draw[very thick](0,-.5)--(0,.6);\draw[decorate,decoration={coil,segment length=5pt,amplitude=3pt}](0,0)--(3,0);\draw[fill=PaleBlue](3,-.4)rectangle(4,.4);\draw[->](3.5,-.9)--(5,-.9)node[right]{$q$};\node at(2,1.2){拉格朗日量中的系数同比增长};
\end{tikzpicture}
\figcap{谐振坐标与指数标度后的振动（题设几何）}\label{advA1:fig:quality-scaled-oscillator-q}
\end{qualityfigure}
\end{advexamplebox}
\begin{advsolution}
\begin{qualityfigure}
\begin{tikzpicture}[>=Latex,font=\small]
\draw[very thick](0,-.5)--(0,.6);\draw[decorate,decoration={coil,segment length=5pt,amplitude=3pt}](0,0)--(3,0);\draw[fill=PaleBlue](3,-.4)rectangle(4,.4);\draw[->](3.5,-.9)--(5,-.9)node[right]{$q$};\node at(2,1.2){拉格朗日量中的系数同比增长};\draw[->](0,-2.4)--(5,-2.4)node[right]{$t$};\draw[DeepBlue,domain=0:5,samples=90]plot(\x,{-2.4+.6*exp(-.3*\x)*cos(240*\x)});\node at(2,-3.3){$q=e^{-\gamma t}\bar q$};
\end{tikzpicture}
\figcap{谐振坐标与指数标度后的振动（解答辅助图）}\label{advA1:fig:quality-scaled-oscillator-a}
\end{qualityfigure}
图~\ref{advA1:fig:quality-scaled-oscillator-a}中，原坐标 $q$ 具有指数包络，标度坐标 $\bar q=e^{\gamma t}q$ 消去了包络；这项坐标变化不能直接把原哈密顿量当成新系统的守恒能量。

(1) 取 $Q=\sqrt m\,q$，则
$L=(\dot Q^2-\omega^2Q^2)/2$。

(2) 方程与哈密顿量为
\[
 \ddot q+2\gamma\dot q+\omega^2q=0,
 \qquad H=\frac m2e^{2\gamma t}(\dot q^2+\omega^2q^2),
\]
\[
 \dot H=-\frac{\partial L}{\partial t}
 =\gamma me^{2\gamma t}(\omega^2q^2-\dot q^2).
\]

(3) 代入 $q=e^{-\gamma t}\bar q$ 后
\[
 L=\frac m2[(\dot{\bar q}-\gamma\bar q)^2-\omega^2\bar q^2]
\]
不显含时间，因而
\[
 \mathcal E=\frac m2[\dot{\bar q}^{\,2}+(\omega^2-\gamma^2)\bar q^2]
 =H+m\gamma e^{2\gamma t}q\dot q
\]
守恒。

(4) 当 $\gamma<\omega$ 时，
\[
 q=e^{-\gamma t}C\cos(\Omega t+\phi),
 \qquad \Omega=\sqrt{\omega^2-\gamma^2}.
\]
代入 $H$ 可见其非常数部分的振幅为 $\mathcal O(\gamma)$，故即使取
$t=1/\gamma$，仍有
$|H(1/\gamma)-H(0)|\le C_1\gamma+C_2\gamma^2\to0$。这是缓变极限，不是对任意有限 $\gamma$ 的精确能量守恒。
还可给出对全部时间成立的界，避免把短时估计直接外推到 $1/\gamma$。
由守恒量 $\mathcal E=m(\dot{\bar q}^{\,2}+\Omega^2\bar q^2)/2$，有
\[
 H-\mathcal E=-m\gamma\bar q\dot{\bar q}+m\gamma^2\bar q^2,
 \quad |H-\mathcal E|\le\frac{\gamma\mathcal E}{\Omega}
 +\frac{2\gamma^2\mathcal E}{\Omega^2}.
\]
因而 $|H(t)-H(0)|\le2\gamma\mathcal E/\Omega+4\gamma^2\mathcal E/\Omega^2$。
在题设固定初值且 $\gamma\to0^+$ 时，$\mathcal E$ 有界、$\Omega\to\omega>0$，
故这个一致界在 $t=1/\gamma$ 时仍趋于零。

\end{advsolution}

\begin{advtheorem}{循环坐标与动量积分}\label{advA1:thm:cyclic}
若 $\partial L/\partial q_f=0$，则 $q_f$ 是循环坐标，且
\[
 p_f=\frac{\partial L}{\partial\dot q_f}=\text{常量}.
\]
反之，若 $p_f$ 对所有运动守恒，则拉格朗日方程给出 $\partial L/\partial q_f=0$。
\end{advtheorem}

\begin{advtheorem}{连续对称性的诺特守恒量}\label{advA1:thm:noether}
设单参数点变换
$\bar q_\alpha=\Lambda_\alpha(q,t,\lambda)$满足
$\Lambda_\alpha(q,t,0)=q_\alpha$，并使拉格朗日量不变，则
\begin{equation}\label{advA1:eq:noether-charge}
 \mathcal Q=\sum_\alpha p_\alpha
 \left.\frac{\partial\Lambda_\alpha}{\partial\lambda}\right|_{\lambda=0}
\end{equation}
是守恒量。
\end{advtheorem}
证明：记
$\delta q_\alpha=(\partial\Lambda_\alpha/\partial\lambda)_0\delta\lambda$。不变性给出
\[
 0=\delta L=\sum_\alpha\left(
 \frac{\partial L}{\partial q_\alpha}\delta q_\alpha
 +p_\alpha\frac{\dd}{\dd t}\delta q_\alpha\right)
 =\frac{\dd}{\dd t}\sum_\alpha p_\alpha\delta q_\alpha,
\]
其中用了 $\dot p_\alpha=\partial L/\partial q_\alpha$。除以常数 $\delta\lambda$ 即得式\eqref{advA1:eq:noether-charge}。
空间平移对应线动量，旋转对应角动量，时间平移对应能量。

\begin{advexamplebox}{对称解绕的旋转圆筒}\label{advA1:ex:unwinding-cylinder}
质量 $M$、半径 $R$ 的薄壁圆筒可绕定轴自由转动。两条相同轻绳对称缠绕，自由端各系质量 $m$ 的小球。初始直绳段长度为零，整体角速度为 $\omega_0$。释放后绳不打滑。（1）求每侧打开长度 $s(t)$；（2）若两侧可打开绳长的总和为 $L$，讨论原题所问“打开后”绳相对固定在筒上的直径的夹角随时间的关系。须说明所选绳端连接模型与公式适用阶段。

\begin{qualityfigure}
\begin{tikzpicture}[>=Latex,font=\small]
\draw[thick](0,0)circle(1.3);\draw[dashed](-1.7,0)--(1.7,0);\draw[very thick](1.3,0)--(1.3,2);\draw[very thick](-1.3,0)--(-1.3,-2);\fill(1.3,2)circle(3pt)node[right]{$m$};\fill(-1.3,-2)circle(3pt)node[left]{$m$};\draw(0,0)--(1.3,0)node[midway,below]{$R$};\node at(0,.4){$M$};\draw[->,Gold](.7,0)arc(0:100:.7)node[midway,above]{$\omega_0$};\node at(3.5,.7){对称解绕；俯视图};
\end{tikzpicture}
\figcap{两侧对称解绕的切线绳段（题设几何）}\label{advA1:fig:quality-unwinding-cylinder-q}
\end{qualityfigure}
\end{advexamplebox}
\AdvHideHeadingtrue
\begin{advsolution}题图~\ref{advA1:fig:quality-unwinding-cylinder-q}中，圆筒两侧的直绳段与筒面相切，另一侧由中心对称得到。绳不打滑联系的是筒转角、切点方位和打开长度，公式只在仍有筒面接触的阶段有效。

\noindent\begin{minipage}{\linewidth}{\bfseries\color{DeepBlue}解答}\par
\begin{center}\begin{minipage}{\linewidth}\centering
\begin{tikzpicture}[scale=.88]
 \draw[very thick] (0,0) circle(1.05);
 \draw[guide] (-1.45,0)--(1.45,0) node[right]{初始固定直径};
 \coordinate (T) at ({1.05*cos(35)},{1.05*sin(35)});
 \draw[curve] (T)--++({2.0*cos(125)},{2.0*sin(125)}) coordinate (P);
 \fill (P) circle(2.5pt) node[above left]{$m$};
 \draw[->,DeepBlue] (.45,0) arc[start angle=0,end angle=35,radius=.45];
 \node at (.68,.19) {$\phi$};
 \draw[<->,TextGray] ($(T)+(.08,.05)$)--($(P)+(.08,.05)$) node[midway,right]{$s$};
 \fill (T) circle(1.8pt) node[right=4pt]{$T$};
 \node at (0,-1.38) {另一侧与图示关于圆心对称};
\end{tikzpicture}
\bookcaption{圆筒上的切点、直绳段和固定直径}\label{advA1:fig:cylinder}
\end{minipage}\end{center}\end{minipage}

如图\ref{advA1:fig:cylinder}，以圆筒转角 $\theta$ 和每侧直绳长 $s$ 为坐标。图中 $\phi$ 是切点半径相对初始固定直径的方位角，直绳切线的方位角为 $\phi+\pi/2$。绳不打滑意味着
$\phi=\theta-s/R$。两球与圆筒的总动能为
\[
 T=\frac12(M+2m)R^2\dot\theta^2
 +ms^2\left(\dot\theta-\frac{\dot s}{R}\right)^2.
\]
$\theta$ 是循环坐标。记
\[
 A=(M+2m)R^2+2ms^2,\quad B=-\frac{2ms^2}{R},
 \quad C=\frac{2ms^2}{R^2},
\]
则 $T=(A\dot\theta^2+2B\dot\theta\dot s+C\dot s^2)/2$。角动量和能量分别为
\[
 J=A\dot\theta+B\dot s=(M+2m)R^2\omega_0,
 \quad E=\frac12(M+2m)R^2\omega_0^2.
\]
消去 $\dot\theta$ 后
\[
 E=\frac12\left(C-\frac{B^2}{A}\right)\dot s^2+\frac{J^2}{2A}.
\]
代入 $A,B,C,J,E$ 化简，得 $\dot s=R\omega_0$，约定 $\omega_0>0$ 并取打开分支，
\[
 \boxed{s(t)=R\omega_0t},\qquad \dot\theta=\omega_0,
 \qquad \phi=\theta-s/R=0.
\]
所以直绳段在惯性系中方向不变。圆筒相对切点半径转过 $\theta-\phi=\omega_0t=s/R$，这是相对初态的转角；直绳与筒上随转直径的有向几何夹角则为
$\beta=\phi+\pi/2-\theta=\pi/2-\omega_0t$（直线夹角按模 $\pi$ 理解）。若每侧可打开的长度为 $L/2$，则打开时间
$t_*=L/(2R\omega_0)$，在这一阶段终点，累计相对转角为 $L/(2R)$，而绳与随转直径的夹角为 $\beta_*=\pi/2-L/(2R)$（模 $\pi$）。

原文没有明确绳端是否固定圆筒，也未交代全部解绕后两球与绳的连接、脱离或绷紧约束。因此以上 $s(t)$ 与相对夹角公式只用于仍在解绕的阶段 $0<t<t_*$；不能不加新条件延伸到 $t>t_*$。若“打开后”指已完全离开筒面，必须先补全绳端约束，再用新的几何约束与守恒量求运动，原给定信息不足以唯一确定该阶段的夹角。
\end{advsolution}

若 $q_f$ 是循环坐标且 $p_f=c_f$，对 $\dot q_f$ 作勒让德变换得罗斯函数
\begin{equation}\label{advA1:eq:routhian}
 \mathcal R=L-p_f\dot q_f.
\end{equation}
用 $p_f=c_f$ 解出 $\dot q_f$并代回后，其余 $f-1$ 个坐标仍满足拉格朗日形式的罗斯方程。

\begin{advexamplebox}{自由转动圆环上小珠的小振动}\label{advA1:ex:routh-hoop}
\emph{题源说明：原资料仅给主题或注明课堂描述；以下条件为依据主题所作的教学补全。}\par
半径 $a$ 的竖直圆环可绕竖直直径无摩擦转动，圆环对该轴的转动惯量为 $I$。质量 $m$ 的小珠在环上滑动。以珠从最低点向上的角位移 $\theta$ 和环的方位角 $\varphi$ 为坐标。在绕竖直轴角动量为 $J$ 时，求最低点附近的振动频率及稳定条件。

\begin{qualityfigure}
\begin{tikzpicture}[>=Latex,font=\small]
\draw[thick](0,0)ellipse(1.5 and 2);\draw[->](0,-2.4)--(0,2.5)node[above]{竖直转轴};\fill(0,0)circle(2pt)node[left]{$O$};\coordinate(P)at(.7949,-1.6961);\draw[dashed](0,0)--(P);\fill(P)circle(3pt)node[right]{$m$};\draw[DeepBlue](0,-.7)arc(-90:-58:.525 and .7);\node at(.34,-.8){$\theta$};\draw[->,Gold](-.7,.7)arc(180:0:.7 and .2)node[right]{$\varphi$};\node[align=left]at(3.3,0){环半径 $a$\\轴惯量 $I$};
\end{tikzpicture}
\figcap{可绕竖直直径转动的圆环与小珠（题设几何）}\label{advA1:fig:quality-routh-hoop-q}
\end{qualityfigure}
\end{advexamplebox}
\AdvHideHeadingtrue
\begin{advsolution}题图~\ref{advA1:fig:quality-routh-hoop-q}中，小珠到竖直轴的距离是 $a\sin\theta$；它使绕轴惯量成为 $I+ma^2\sin^2\theta$，消去循环坐标时应固定角动量 $J$。

\noindent\begin{minipage}{\linewidth}{\bfseries\color{DeepBlue}解答}\par
\begin{center}\begin{minipage}{\linewidth}\centering
\begin{tikzpicture}[scale=.9]
 \draw[axis] (0,2)--(0,-2.1) node[below]{转轴};
 \draw[very thick] (0,0) ellipse (1.35 and 1.8);
 \coordinate (P) at ({1.35*sin(32)},{-1.8*cos(32)});
 \draw[guide] (0,0)--(P);\fill (0,0) circle(1.5pt) node[left]{$O$};
 \fill (P) circle(2.4pt) node[right=5pt]{$m$};
 \draw[DeepBlue,->] (0,-.68) arc[start angle=-90,end angle=-58,x radius=.51,y radius=.68];
 \node at (.36,-.82) {$\theta$};
 \draw[TextGray,->] ({.75*cos(160)},{.65+.20*sin(160)}) arc[start angle=160,end angle=20,x radius=.75,y radius=.20];
 \node at (.90,.77) {$\varphi$};
\end{tikzpicture}
\bookcaption{可绕竖直直径转动的圆环与环上小珠}\label{advA1:fig:hoop}
\end{minipage}\end{center}\end{minipage}

如图\ref{advA1:fig:hoop}，
\[
 L=\frac12ma^2\dot\theta^2
 +\frac12(I+ma^2\sin^2\theta)\dot\varphi^2
 +mga\cos\theta .
\]
$\varphi$ 循环，
$J=(I+ma^2\sin^2\theta)\dot\varphi$。作罗斯约化：
\[
 \mathcal R=\frac12ma^2\dot\theta^2+mga\cos\theta
 -\frac{J^2}{2(I+ma^2\sin^2\theta)}.
\]
在 $\theta=0$ 附近保留二阶项，
\[
 \mathcal R=\text{常量}+\frac12ma^2\dot\theta^2
 -\frac12ma^2\left(\frac ga-\frac{J^2}{I^2}\right)\theta^2+\order{\theta^4}.
\]
因而
\[
 \boxed{\omega^2=\frac ga-\left(\frac JI\right)^2},
 \qquad \frac{|J|}I<\sqrt{\frac ga}.
\]
当 $|J|/I=\sqrt{g/a}$ 时二阶恢复力消失，但有效势的首个非零增量为
\[
 U_{\rm eff}(\theta)-U_{\rm eff}(0)
 =mga\left(\frac18+\frac{ma^2}{2I}\right)\theta^4+\order{\theta^6}>0
\]
（充分小的非零 $\theta$），所以最低点仍为稳定平衡，只是线性小振动频率为零，有限小振幅的周期须由非线性项决定。因此稳定范围含等号；上式严格不等号对应非零的线性回复频率。超过阈值时最低点变为不稳定点。
\end{advsolution}

\subsection{小振动、傅里叶变换与周期系统}

在稳定平衡点 $q_\alpha=0$ 附近，将动能和势能展开到二阶：
\begin{equation}\label{advA1:eq:small-L}
 L=\frac12\dot{\vect q}^{\mathsf T}\mathsf T\dot{\vect q}
 -\frac12\vect q^{\mathsf T}\mathsf V\vect q .
\end{equation}
$\mathsf T$ 是正定的惯性矩阵，$\mathsf V$ 是对称的弹性矩阵。被忽略的项至少含有三个小量，可用微扰迭代逐阶修正。

\begin{advdefinition}{线性变换}\label{advA1:def:linear-map}
向量空间上满足
$F(\lambda\vect x+\mu\vect y)=\lambda F(\vect x)+\mu F(\vect y)$
的映射称为线性变换。
\end{advdefinition}

\begin{advtheorem}{线性变换的矩阵表示}\label{advA1:thm:linear-matrix}
在选定基底后，每个有限维线性变换唯一对应一个矩阵 $\mathsf F$，使
$\vect y=F(\vect x)$ 等价于 $[\vect y]=\mathsf F[\vect x]$。
\end{advtheorem}
证明：只需把每个基矢的像作为矩阵的一列；线性性保证它对任意向量成立，基底展开的唯一性保证矩阵唯一。

\begin{advdefinition}{对称二次型}\label{advA1:def:quadratic-form}
若 $S(\vect x,\vect y)$ 对两个变量均线性，且
$S(\vect x,\vect y)=S(\vect y,\vect x)$，则称其为对称二次型对应的双线性型。
\end{advdefinition}

\begin{advtheorem}{二次型的对称矩阵表示}\label{advA1:thm:quadratic-matrix}
在选定基底后，存在唯一对称矩阵 $\mathsf S$，使
\[
 S(\vect x,\vect y)=\vect x^{\mathsf T}\mathsf S\vect y.
\]
\end{advtheorem}
证明：令 $S_{ij}=S(\vect e_i,\vect e_j)$，用双线性展开即得矩阵表示；对称性给出 $S_{ij}=S_{ji}$。

\begin{advdefinition}{合同变换}\label{advA1:def:congruence}
对非奇异矩阵 $\mathsf P$，
\[
 \mathsf S\longmapsto \mathsf P^{\mathsf T}\mathsf S\mathsf P
\]
称为对称矩阵的合同变换。它是二次型在坐标变换 $\vect q=\mathsf P\vect Q$ 下的自然变换。
\end{advdefinition}

\begin{advtheorem}{惯性矩阵与弹性矩阵的同时对角化}\label{advA1:thm:normal-modes}
若 $\mathsf T$ 对称正定、$\mathsf V$ 对称，则存在非奇异矩阵 $\mathsf A$，使
\begin{equation}\label{advA1:eq:sim-diag}
 \mathsf A^{\mathsf T}\mathsf T\mathsf A=\mathsf I,
 \qquad
 \mathsf A^{\mathsf T}\mathsf V\mathsf A
 =\operatorname{Diag}(\omega_1^2,\ldots,\omega_n^2).
\end{equation}
\end{advtheorem}
证明：因 $\mathsf T$ 正定，存在可逆平方根 $\mathsf T^{1/2}$。矩阵
$\mathsf K=\mathsf T^{-1/2}\mathsf V\mathsf T^{-1/2}$ 对称，因此可被正交矩阵 $\mathsf O$ 对角化。取
$\mathsf A=\mathsf T^{-1/2}\mathsf O$，立即得式\eqref{advA1:eq:sim-diag}。令 $\vect q=\mathsf A\vect Q$，式\eqref{advA1:eq:small-L}化为
\[
 L=\frac12\sum_j(\dot Q_j^2-\omega_j^2Q_j^2).
\]
$\omega_j^2>0$ 表示稳定振动，$\omega_j^2=0$ 表示零模，$\omega_j^2<0$ 表示指数不稳定。

\begin{advexamplebox}{刚体复摆的小振动}\label{advA1:ex:physical-pendulum}
\emph{题源说明：原资料仅给主题或注明课堂描述；以下条件为依据主题所作的教学补全。}\par
刚体质量为 $M$，质心到水平固定轴的距离为 $a$，对该轴的转动惯量为 $I$。求竖直向下稳定位置附近的周期。

\begin{qualityfigure}
\begin{tikzpicture}[>=Latex,font=\small]
\fill(0,2.5)circle(2pt)node[above]{$O$};\draw[dashed](0,2.5)--(0,-1);\coordinate(P)at(1.5,.1);\draw[thick](0,2.5)--(P);\fill(P)circle(3pt)node[right]{$G$};\draw(0,1.6)arc(-90:-58:.9)node[midway,below right]{$\theta$};\node at(3,1.2){$OG=a$};\node at(3,.4){轴惯量 $I$，质量 $M$};
\end{tikzpicture}
\figcap{固定轴与偏离铅垂线的刚体质心（题设几何）}\label{advA1:fig:quality-physical-pendulum-q}
\end{qualityfigure}
\end{advexamplebox}
\begin{advsolution}
\begin{qualityfigure}
\begin{tikzpicture}[>=Latex,font=\small]
\fill(0,2.5)circle(2pt)node[above]{$O$};\draw[dashed](0,2.5)--(0,-1);\coordinate(P)at(1.5,.1);\draw[thick](0,2.5)--(P);\fill(P)circle(3pt)node[right]{$G$};\draw(0,1.6)arc(-90:-58:.9)node[midway,below right]{$\theta$};\node at(3,1.2){$OG=a$};\node at(3,.4){轴惯量 $I$，质量 $M$};\draw[->,Gold](P)--++(0,-1)node[below]{$Mg$};
\end{tikzpicture}
\figcap{固定轴与偏离铅垂线的刚体质心（解答辅助图）}\label{advA1:fig:quality-physical-pendulum-a}
\end{qualityfigure}
图~\ref{advA1:fig:quality-physical-pendulum-a}中，质心 $G$ 距支点轴为 $a$，重力矩的力臂是 $a\sin\theta$；小角展开后用刚体对该固定轴的惯量 $I$，不能用质点式 $Ma^2$ 代替。

以偏角 $\theta$ 为坐标，
\[
 L=\frac12I\dot\theta^2+Mga\cos\theta.
\]
精确方程是 $I\ddot\theta+Mga\sin\theta=0$。小角度下
\[
 I\ddot\theta+Mga\theta=0,
 \qquad \boxed{\omega=\sqrt{\frac{Mga}{I}}},
 \qquad \boxed{T_0=2\pi\sqrt{\frac I{Mga}}}.
\]
若 $I=Ma^2$，回到长度 $a$ 的单摆；若 $a\to0$，重力恢复力矩消失，周期趋于无穷，两个极限均正确。
\end{advsolution}

\subsubsection{陀螺型速度耦合与傅科摆}

考虑二维谐振子
\[
 T=\frac m2(\dot x^2+\dot y^2),\qquad
 V_0=\frac k2(x^2+y^2),
\]
并受到速度型力
\[
 F_x=-m\Omega\dot y,\qquad F_y=m\Omega\dot x.
\]
它可由含速势
\[
 V_g=\frac{m\Omega}{2}(x\dot y-y\dot x)
\]
描述。令 $z=x+\ii y$，方程化为
\begin{equation}\label{advA1:eq:foucault-complex}
 \ddot z-\ii\Omega\dot z+\omega^2z=0,
 \qquad \omega^2=k/m.
\end{equation}
作变换 $z=e^{\ii\Omega t/2}\zeta$，得
\[
 \ddot\zeta+\omega_0^2\zeta=0,
 \qquad \omega_0^2=\omega^2+\frac{\Omega^2}{4}.
\]
两个圆偏振模的频率为 $\omega_0\pm\Omega/2$，线偏振叠加后的振动平面以 $\Omega/2$ 进动。对地球上的傅科摆，取 $x$ 向东、$y$ 向北、竖直向上为正，地球自转的竖直分量为 $\Omega_E\sin\lambda$。科里奥利力的水平分量为 $F_x=2m\Omega_E\sin\lambda\,\dot y$、$F_y=-2m\Omega_E\sin\lambda\,\dot x$，故按本节的力和复坐标约定，应代入
$\Omega=-2\Omega_E\sin\lambda$，振动平面的有符号进动角速度为 $-\Omega_E\sin\lambda$；北半球从上方看为顺时针，进动速率的大小为 $|\Omega_E\sin\lambda|$。

\begin{advtheorem}{周期函数的傅里叶级数}\label{advA1:thm:fourier-series}
对周期 $l$ 的光滑函数，
\[
 f(x)=\sum_{n=-\infty}^{\infty}c_n e^{\ii2\pi nx/l},
 \qquad c_n=\frac1l\int_0^l f(x)e^{-\ii2\pi nx/l}\dd x.
\]
\end{advtheorem}
证明要点是函数 $e^{\ii2\pi nx/l}$ 在 $[0,l]$ 上正交：
\[
 \frac1l\int_0^l e^{\ii2\pi(n-m)x/l}\dd x=\delta_{nm}.
\]
将展开式两边与第 $m$ 个基函数的共轭相乘并积分，即得 $c_m$。光滑性保证系数足够快衰减并使级数按所需意义收敛。

\begin{advtheorem}{平方可积函数的傅里叶变换}\label{advA1:thm:fourier-transform}
采用对称归一化，
\[
 g(k)=\frac1{\sqrt{2\pi}}\int_{-\infty}^{\infty}f(x)e^{-\ii kx}\dd x,
 \qquad
 f(x)=\frac1{\sqrt{2\pi}}\int_{-\infty}^{\infty}g(k)e^{\ii kx}\dd k.
\]
对 $f\in L^2(\mathbb R)$，等式按 $L^2$ 意义成立。
\end{advtheorem}
可先在长区间 $[-L/2,L/2]$ 上作傅里叶级数，再令
$L\to\infty$。此时离散波数间距 $\Delta k=2\pi/L\to0$，级数的黎曼和趋于积分，得到上式。更严格的表述依赖平方可积函数空间的完备性。

连续/离散、周期/非周期的四种组合分别对应傅里叶变换、傅里叶级数、离散时间傅里叶变换和离散傅里叶变换。它们的共同核心都是用平移对称性的本征函数实现解耦。

\subsubsection{一维单原子链}

在平衡位置 $x_n^{(0)}=nl$ 上放置质量 $m$ 的质点，纵向位移为 $u_n$，最近邻弹簧常量为 $\kappa$。
\[
 L=\sum_n\left[\frac12m\dot u_n^2-\frac12\kappa(u_{n+1}-u_n)^2\right].
\]
对离散坐标作变换
\[
 \phi(k,t)=\sum_{n=-\infty}^{\infty}u_n(t)e^{-\ii knl},
 \quad
 u_n(t)=\frac l{2\pi}\int_{-\pi/l}^{\pi/l}
 \phi(k,t)e^{\ii knl}\dd k.
\]
由于 $\phi(k+2\pi/l)=\phi(k)$，只需保留第一布里渊区 $[-\pi/l,\pi/l]$。代入并利用正交性，各个 $k$ 模彼此解耦，色散关系为
\begin{equation}\label{advA1:eq:lattice-dispersion}
 \boxed{\omega(k)=2\sqrt{\frac\kappa m}
 \left|\sin\frac{kl}{2}\right|}.
\end{equation}
长波极限 $|kl|\ll1$ 下，$\omega\simeq l\sqrt{\kappa/m}|k|$，群速度趋于
$l\sqrt{\kappa/m}$；在布里渊区边界 $k=\pm\pi/l$ 处群速度为零。

\begin{advexamplebox}{用极坐标求傅科摆进动}\label{advA1:ex:foucault-polar}
对 $L=m(\dot x^2+\dot y^2)/2-m\omega^2(x^2+y^2)/2-m\Omega(x\dot y-y\dot x)/2$，改用极坐标 $(r,\theta)$，求共转线偏振解的进动速率，并说明一般初态及静止释放时的区别。

\begin{qualityfigure}
\begin{tikzpicture}[>=Latex,font=\small]
\draw[->](-2.5,0)--(3,0)node[right]{$x$};\draw[->](0,-2)--(0,2.4)node[above]{$y$};\draw[DeepBlue,rotate=25](0,0)ellipse(2 and .5);\draw[dashed,rotate=25](-2.3,0)--(2.3,0)node[right]{共转长轴};\draw(.8,0)arc(0:25:.8)node[midway,right]{$\theta$};
\end{tikzpicture}
\figcap{共转坐标中椭圆的主轴与极角（题设几何）}\label{advA1:fig:quality-foucault-polar-q}
\end{qualityfigure}
\end{advexamplebox}
\begin{advsolution}
\begin{qualityfigure}
\begin{tikzpicture}[>=Latex,font=\small]
\draw[->](-2.5,0)--(3,0)node[right]{$x$};\draw[->](0,-2)--(0,2.4)node[above]{$y$};\draw[DeepBlue,rotate=25](0,0)ellipse(2 and .5);\draw[dashed,rotate=25](-2.3,0)--(2.3,0)node[right]{共转长轴};\draw(.8,0)arc(0:25:.8)node[midway,right]{$\theta$};\node[align=center]at(0,-2.6){共转椭圆的主轴以 $\Omega/2$ 转动\\一般初态不必是直线振动};
\end{tikzpicture}
\figcap{共转坐标中椭圆的主轴与极角（解答辅助图）}\label{advA1:fig:quality-foucault-polar-a}
\end{qualityfigure}
图~\ref{advA1:fig:quality-foucault-polar-a}中，椭圆的轴角与质点的瞬时极角不同；只有共转系中的特定线偏振解才有 $J=0$，一般初态应保留椭圆的横向运动。

因 $x\dot y-y\dot x=r^2\dot\theta$，
\[
 L=\frac m2(\dot r^2+r^2\dot\theta^2)-\frac k2r^2
 -\frac{m\Omega}{2}r^2\dot\theta,
 \qquad J=mr^2\left(\dot\theta-\frac\Omega2\right).
\]
令共转角 $\psi=\theta-\Omega t/2$，则 $J=mr^2\dot\psi$。
在共转坐标中选取线偏振振动才有 $J=0$；除穿过原点时极角的 $\pi$ 跳变外，
振动直线满足 $\dot\theta=\Omega/2$。
此时用可正可负的直线坐标 $q$ 代替非负半径，得到
\[
 \ddot q+\omega_0^2q=0,\qquad \omega_0^2=\omega^2+\Omega^2/4.
\]
一般初态在共转系中形成固定椭圆，该椭圆的轴随坐标系以 $\Omega/2$ 转动。
例如原坐标中从 $z(0)=a,\dot z(0)=0$ 静止释放，有
\[
 J=-\frac12m\Omega a^2\ne0,\qquad
 z(t)=a\e^{\ii\Omega t/2}
 \left[\cos\omega_0t-\frac{\ii\Omega}{2\omega_0}\sin\omega_0t\right].
\]
在 $|\Omega|\ll\omega_0$ 时椭圆很窄，才可近似称为单一振动平面进动。
因此 $J=0$ 是特定线偏振解的条件，不能由任意“静止释放”推出。

\end{advsolution}

\begin{advexamplebox}{无限一维电阻网络}\label{advA1:ex:resistor-chain}
\emph{题源说明：原资料仅给主题或注明课堂描述；以下条件为依据主题所作的教学补全。}\par
无限多个节点按整数 $n$ 排列，相邻节点之间接电阻 $R$，每个节点再通过电阻 $r$ 接公共地。求 $n=0$ 节点对地的等效电阻。

\begin{qualityfigure}
\begin{tikzpicture}[>=Latex,font=\small]
\foreach\i in{-2,-1,0,1,2}{\fill(1.5*\i,0)circle(2pt)node[above]{$\i$};\draw(1.5*\i,0)--(1.5*\i,-.3)--++(.12,-.12)--++(-.24,-.2)--++(.24,-.2)--++(-.24,-.2)--++(.12,-.12)--(1.5*\i,-1.5);\node[right]at(1.5*\i,-.8){$r$};}\draw(-3,-1.5)--(3,-1.5)node[right]{公共地};\foreach\i in{-2,-1,0,1}{\draw(1.5*\i,0)--++(.4,0)--++(.1,.15)--++(.2,-.3)--++(.2,.3)--++(.2,-.3)--++(.1,.15)--++(.3,0);\node at(1.5*\i+.75,.5){$R$};}\node at(-3.7,0){$\cdots$};\node at(3.7,0){$\cdots$};
\end{tikzpicture}
\figcap{无限电阻链的节点和对地支路（题设几何）}\label{advA1:fig:quality-resistor-chain-q}
\end{qualityfigure}
\end{advexamplebox}
\begin{advsolution}
\begin{qualityfigure}
\begin{tikzpicture}[>=Latex,font=\small]
\foreach\i in{-2,-1,0,1,2}{\fill(1.5*\i,0)circle(2pt)node[above]{$\i$};\draw(1.5*\i,0)--(1.5*\i,-.3)--++(.12,-.12)--++(-.24,-.2)--++(.24,-.2)--++(-.24,-.2)--++(.12,-.12)--(1.5*\i,-1.5);\node[right]at(1.5*\i,-.8){$r$};}\draw(-3,-1.5)--(3,-1.5)node[right]{公共地};\foreach\i in{-2,-1,0,1}{\draw(1.5*\i,0)--++(.4,0)--++(.1,.15)--++(.2,-.3)--++(.2,.3)--++(.2,-.3)--++(.1,.15)--++(.3,0);\node at(1.5*\i+.75,.5){$R$};}\node at(-3.7,0){$\cdots$};\node at(3.7,0){$\cdots$};\draw[Gold,thick](-.3,-1.7)rectangle(.3,.7);\node at(0,-2.3){在 $n=0$ 注入电流，两侧衰减解须在无穷远有界};
\end{tikzpicture}
\figcap{无限电阻链的节点和对地支路（解答辅助图）}\label{advA1:fig:quality-resistor-chain-a}
\end{qualityfigure}
图~\ref{advA1:fig:quality-resistor-chain-a}中，每个整数节点均通过 $r$ 接地，相邻节点通过 $R$ 相连；在零号节点注入电流后，两侧节点电势必须满足同一差分方程及无穷远有界条件。

向 $0$ 节点注入电流 $I$，节点电势为 $V_n$。对所有 $n$，
\[
 \frac{V_n}{r}+\frac{2V_n-V_{n-1}-V_{n+1}}R=I\delta_{n0}.
\]
作离散傅里叶变换，
\[
 V(k)=\frac I{r^{-1}+2R^{-1}(1-\cos k)}.
\]
因此
\[
 \frac{V_0}{I}=\frac1{2\pi}\int_{-\pi}^{\pi}
 \frac{\dd k}{r^{-1}+2R^{-1}(1-\cos k)}.
\]
% A1CompletenessInput20261024: standard-integral
% A1Completeness20261024: resistor-integral
下面把最后的积分算清，避免把电路答案藏在“标准积分”中。对$a>|b|$，设
\[
 J(a,b)=\frac1{2\pi}\int_{-\pi}^{\pi}\frac{\dd k}{a-b\cos k}.
\]
令$t=\tan(k/2)$，则$\cos k=(1-t^2)/(1+t^2)$、$\dd k=2\dd t/(1+t^2)$。
$k$从$-\pi$走到$\pi$时，$t$从$-\infty$走到$+\infty$，故
\[
 J=\frac1\pi\int_{-\infty}^{\infty}
 \frac{\dd t}{(a-b)+(a+b)t^2}.
\]
因$a-b,a+b$都为正，再令$u=t\sqrt{(a+b)/(a-b)}$：
\[
 J=\frac1{\pi\sqrt{a^2-b^2}}
 \int_{-\infty}^{\infty}\frac{\dd u}{1+u^2}
 =\frac{\arctan u|_{-\infty}^{\infty}}{\pi\sqrt{a^2-b^2}}
 =\frac1{\sqrt{a^2-b^2}}.
\]
\begin{qualityfigure}
\begin{tikzpicture}[>=Latex,font=\small]
 \begin{scope}
  \draw[->](-2.2,1)--(2.3,1)node[right]{$k$};
  \foreach\x/\lab in{-1.8/-\pi,0/0,1.8/\pi}{\draw(\x,.92)--(\x,1.08);\node[above]at(\x,1.1){$\lab$};}
  \draw[->](-2.2,-.7)--(2.3,-.7)node[right]{$t$};
  \node[below]at(-1.8,-.8){$-\infty$};\node[below]at(0,-.8){$0$};\node[below]at(1.8,-.8){$+\infty$};
  \foreach\x in{-1.8,0,1.8}{\draw[->,DeepBlue](\x,.65)--(\x,-.35);}
  \node at(0,-1.7){$t=\tan(k/2)$};
 \end{scope}
 \begin{scope}[shift={(6.2,-.4)}]
  \draw[->](-2.1,0)--(2.3,0)node[right]{$u$};
  \draw[->](0,0)--(0,1.65)node[above]{$(1+u^2)^{-1}$};
  \fill[PaleBlue,domain=-2:2,samples=80](-2,0)--plot(\x,{1/(1+\x*\x)})--(2,0)--cycle;
  \draw[DeepBlue,thick,domain=-2:2,samples=80]plot(\x,{1/(1+\x*\x)});
  \node at(0,-.65){全实轴曲线下的面积为$\pi$};
 \end{scope}
\end{tikzpicture}
\figcap{半角代换将一个周期的积分化为全实轴上的反正切积分；右图的两端尾部延伸到无穷远。}\label{advA1:fig:completeness-resistor-integral}
\end{qualityfigure}
图~\ref{advA1:fig:completeness-resistor-integral}左侧的三个对应点固定了积分区间，右侧面积由反正切的两个端点极限得到。本题取$a=r^{-1}+2R^{-1}$、$b=2R^{-1}$，电阻为正保证$a-b=r^{-1}>0$，代入得


\[
 \boxed{R_{\rm eq}=\frac1{\sqrt{r^{-2}+4/(rR)}}
 =\frac r{\sqrt{1+4r/R}}}.
\]
$R\to\infty$ 时各节点断开，$R_{\rm eq}\to r$；$r\to\infty$ 时整个网络没有到地通路，$R_{\rm eq}\to\infty$，极限检查通过。
\end{advsolution}

\begin{advexamplebox}{六原子环的简正模}\label{advA1:ex:six-atom-ring}
\emph{题源说明：原资料仅给主题或注明课堂描述；以下条件为依据主题所作的教学补全。}\par
六个质量均为 $m$ 的原子排成闭环，只考虑沿环方向的小位移 $u_j$，相邻原子之间的有效弹簧常量为 $\kappa$，$u_{j+6}=u_j$。求全部简正频率及简并度。

\begin{qualityfigure}
\begin{tikzpicture}[>=Latex,font=\small]
\foreach\i in{0,...,5}{\coordinate(P\i)at({60*\i}:1.8);\fill(P\i)circle(2.5pt)node[below right]{$j=\i$};\draw[->,Gold](P\i)--++({60*\i+90}:.55)node[above]{$u_\i$};}\draw[DeepBlue,thick](P0)--(P1)--(P2)--(P3)--(P4)--(P5)--cycle;\node[align=center]at(0,0){$m_j=m$\\$\kappa$};\node at(0,-2.7){沿环方向的小位移，$u_{j+6}=u_j$};
\end{tikzpicture}
\figcap{六原子的周期边界与切向小位移（题设几何）}\label{advA1:fig:quality-six-atom-ring-q}
\end{qualityfigure}
\end{advexamplebox}
\AdvHideHeadingtrue
\begin{advsolution}题图~\ref{advA1:fig:quality-six-atom-ring-q}中，相邻原子之间的力取决于位移差，首尾相接给出 $u_{j+6}=u_j$；允许波数由这个周期边界量子化，不能使用自由端链的边界条件。

\noindent\begin{minipage}{\linewidth}{\bfseries\color{DeepBlue}解答}\par
\begin{center}\begin{minipage}{\linewidth}\centering
\begin{tikzpicture}[scale=.92]
 \foreach \j in {0,...,5}{
   \coordinate (P\j) at ({90-60*\j}:1.42);
 }
 \foreach \j [evaluate=\j as \k using {int(mod(\j+1,6))}] in {0,...,5}{
   \draw[spring] (P\j)--(P\k);
 }
 \foreach \j in {0,...,5}{
   \fill (P\j) circle(3pt);
   \node at ($1.28*(P\j)$) {$u_{\j}$};
 }
 \draw[->,DeepBlue] (0,2.15) arc[start angle=90,end angle=25,radius=2.15];
 \node[above right] at (2.15,1.85) {正向};
\end{tikzpicture}
\bookcaption{六原子环中的周期边界与最近邻耦合}\label{advA1:fig:ring}
\end{minipage}\end{center}\end{minipage}

如图\ref{advA1:fig:ring}，方程为
\[
 m\ddot u_j=\kappa(u_{j+1}+u_{j-1}-2u_j).
\]
取 $u_j=Ae^{\ii(kj-\omega t)}$，周期边界给出
$k_n=2\pi n/6$，$n=0,\ldots,5$，而
\[
 \omega_n=2\sqrt{\frac\kappa m}\left|\sin\frac{k_n}{2}\right|.
\]
所以频谱为
\[
 \boxed{0\ (1),\quad \sqrt{\kappa/m}\ (2),\quad
 \sqrt{3\kappa/m}\ (2),\quad 2\sqrt{\kappa/m}\ (1)}.
\]
括号中是简并度。零模对应六个原子沿环同向等量移动，弹簧不变形；$n$ 与 $6-n$ 构成反向传播的一对简并模。
\end{advsolution}

\section{刚体}

\subsection{质量几何与矩生成函数}

对密度分布 $\rho(x,y,z)$，定义笛卡尔矩
\begin{equation}\label{advA1:eq:cartesian-moment}
 M_{pqr}=\iiint \rho(x,y,z)x^py^qz^r\dd V .
\end{equation}
零阶矩是总质量，一阶矩确定质心，二阶矩决定转动惯量张量。更高阶矩则继续编码质量分布的细节。

\begin{advexamplebox}{矩的平移定理}\label{advA1:ex:moment-translation}
设相对原点 $O$ 的矩为 $M_{pqr}$。把物体整体平移 $\vect a=(a,b,c)$ 后，证明新矩为
\[
 M'_{pqr}=\sum_{i=0}^{p}\sum_{j=0}^{q}\sum_{k=0}^{r}
 \binom pi\binom qj\binom rk
 a^{p-i}b^{q-j}c^{r-k}M_{ijk}.
\]
\end{advexamplebox}
\begin{advsolution}
平移后同一质量元的坐标由 $(x,y,z)$ 变为 $(x+a,y+b,z+c)$，而 $\rho\dd V$ 不变。代入定义\eqref{advA1:eq:cartesian-moment}，对三个因子分别用二项式定理：
\[
 (x+a)^p(y+b)^q(z+c)^r
 =\sum_{ijk}\binom pi\binom qj\binom rk
 x^iy^jz^k a^{p-i}b^{q-j}c^{r-k}.
\]
逐项积分即得所证公式。对一阶矩，它给出
$\vect M'_1=\vect M_1+M\vect a$；对二阶矩，它给出熟悉的平行轴项。
\end{advsolution}

引入矩生成函数
\begin{equation}\label{advA1:eq:moment-generator}
 P(s_1,s_2,s_3)=\iiint\rho(x,y,z)
 e^{s_1x+s_2y+s_3z}\dd V .
\end{equation}

\begin{advtheorem}{笛卡尔矩的生成}\label{advA1:thm:moment-generator}
\[
 M_{pqr}=\left.
 \frac{\partial^{p+q+r}P}
 {\partial s_1^p\partial s_2^q\partial s_3^r}
 \right|_{(0,0,0)}.
\]
\end{advtheorem}
证明：在积分号内对 $s_1,s_2,s_3$ 求导，每次分别带出一个 $x,y,z$；最后令 $s_1=s_2=s_3=0$，指数因子变为 $1$，正好得到式\eqref{advA1:eq:cartesian-moment}。因而在收敛条件满足时，全部矩就是 $P$ 在原点处的多元泰勒系数。

\begin{advexamplebox}{均匀正方形薄板的前四级矩}\label{advA1:ex:square-moments}
在 $0\le x\le l, 0\le y\le l$ 上面密度 $\sigma$ 均匀分布。求总阶数 $p+q=0,1,2,3$ 的全部平面笛卡尔矩 $M_{pq}$。
\end{advexamplebox}
\begin{advsolution}
\[
 M_{pq}=\sigma\int_0^l x^p\dd x\int_0^l y^q\dd y
 =\frac{\sigma l^{p+q+2}}{(p+1)(q+1)}.
\]
因此
\begin{align*}
 &M_{00}=\sigma l^2,\\
 &M_{10}=M_{01}=\frac12\sigma l^3,\\
 &M_{20}=M_{02}=\frac13\sigma l^4,
 \qquad M_{11}=\frac14\sigma l^4,\\
 &M_{30}=M_{03}=\frac14\sigma l^5,
 \qquad M_{21}=M_{12}=\frac16\sigma l^5.
\end{align*}
单位为质量乘长度的 $p+q$ 次方，与定义相符。
\end{advsolution}

\begin{advexamplebox}{卷积与矩生成函数}\label{advA1:ex:convolution}
定义三维卷积
\[
 \gamma(\vect r)=(\alpha*\beta)(\vect r)
 =\int\alpha(\vect r')\beta(\vect r-\vect r\,')\dd^3r'.
\]
若三个分布的矩生成函数分别为 $P,Q,R$，矩分别为 $L,M,N$，证明 $R=PQ$，并写出 $N_{pqr}$。
\end{advexamplebox}
\begin{advsolution}
对卷积作变量代换 $\vect u=\vect r-\vect r\,'$，
\begin{align*}
 R(\vect s)
 &=\int\dd^3r\int\dd^3r'\,\alpha(\vect r')
 \beta(\vect r-\vect r\,')e^{\vect s\cdot\vect r}\\
 &=\left[\int\alpha(\vect r')e^{\vect s\cdot\vect r'}\dd^3r'\right]
 \left[\int\beta(\vect u)e^{\vect s\cdot\vect u}\dd^3u\right]
 =P(\vect s)Q(\vect s).
\end{align*}
对乘积求 $p+q+r$ 阶偏导并用多重莱布尼茨公式：
\[
 N_{pqr}=\sum_{i=0}^{p}\sum_{j=0}^{q}\sum_{k=0}^{r}
 \binom pi\binom qj\binom rk
 L_{ijk}M_{p-i,q-j,r-k}.
\]
这一关系与“独立随机变量之和的矩生成函数相乘”是同一个结构。
\end{advsolution}

\subsection{欧拉运动学方程}

对定点转动，取初始时刻的固连正交基
$(\vect e_1,\vect e_2,\vect e_3)$。设形变映射为 $\vect R(\vect r,t)$，并定义
$\vect\varepsilon_i=\partial\vect R/\partial x_i$。

\begin{advdefinition}{刚体}\label{advA1:def:rigid-body}
刚体的形变映射在每一时刻都保长且保向：
\[
 \vect\varepsilon_i\cdot\vect\varepsilon_j=\delta_{ij},
 \qquad (\vect\varepsilon_1\times\vect\varepsilon_2)
 \cdot\vect\varepsilon_3=+1.
\]
\end{advdefinition}
固定原点后，这种映射由特殊正交矩阵 $\mathsf A(t)\in SO(3)$ 描述：
\[
 \vect R=\mathsf A(t)\vect r,
 \qquad \mathsf A^{\mathsf T}\mathsf A=\mathsf I,
 \qquad \det\mathsf A=1.
\]
“保长”给出正交性，“保向”排除行列式为 $-1$ 的镜像变换。连续的刚体运动不能穿过 $\det\mathsf A=0$，因而也不能在运动中突然改变左右手性。

\begin{advexamplebox}{三条坐标轴的转动矩阵}\label{advA1:ex:axis-rotations}
写出绕 $x,y,z$ 轴主动旋转 $\theta$ 的矩阵，检验不同轴转动是否对易，并讨论能否用一组实坐标同时将所有转动对角化。
\end{advexamplebox}
\begin{advsolution}
\[
 R_x=\begin{pmatrix}1&0&0\\0&c&-s\\0&s&c\end{pmatrix},
 \quad
 R_y=\begin{pmatrix}c&0&s\\0&1&0\\-s&0&c\end{pmatrix},
 \quad
 R_z=\begin{pmatrix}c&-s&0\\s&c&0\\0&0&1\end{pmatrix},
\]
其中 $c=\cos\theta, s=\sin\theta$。例如对矢量 $\vect e_x$，
$R_xR_y\vect e_x\ne R_yR_x\vect e_x$，所以一般 $R_xR_y\ne R_yR_x$。无穷小角时对易子是二阶小量，这也解释了小转角下“转角矢量近似可相加”只是一阶结论。

一般实转动的本征值为 $1,e^{\pm\ii\theta}$，除 $\theta=0,\pi$ 等特例外不能在实数域对角化。更不可能用同一组基同时对角化全部 $SO(3)$ 元素，因为可同时对角化的矩阵必彼此对易。
\end{advsolution}

定义刚体坐标中的角速度矩阵
\begin{equation}\label{advA1:eq:omega-matrix}
 \bm{\Omega}=\mathsf{A}^{\mathsf T}\dot{\mathsf{A}}.
\end{equation}

\begin{advtheorem}{$SO(3)$ 的生成元是反对称矩阵}\label{advA1:thm:so3-generator}
\[
 \bm{\Omega}^{\mathsf T}=-\bm{\Omega}.
\]
\end{advtheorem}
对 $\mathsf A^{\mathsf T}\mathsf A=\mathsf I$ 求导：
$\dot{\mathsf A}^{\mathsf T}\mathsf A+\mathsf A^{\mathsf T}\dot{\mathsf A}=0$。第一项正是第二项的转置，因而得证。三维反对称矩阵只有三个独立分量，可与伪矢量 $\vect\omega=(\omega_1,\omega_2,\omega_3)$ 对应：
\[
 [\vect\omega]_\times=
 \begin{pmatrix}
 0&-\omega_3&\omega_2\\
 \omega_3&0&-\omega_1\\
 -\omega_2&\omega_1&0
 \end{pmatrix},
 \qquad [\vect\omega]_\times\vect r=\vect\omega\times\vect r.
\]
所以刚体上的点在刚体坐标中满足 $\vect v=\vect\omega\times\vect r$。

\begin{advexamplebox}{矩阵指数与有限转动}\label{advA1:ex:matrix-exponential}
记 $[\vect n]_\times$ 为单位矢量 $\vect n$ 的叉乘矩阵。证明
\[
 R_{\vect n}(\theta)=e^{\theta[\vect n]_\times}
 =\lim_{N\to\infty}\left(\mathsf I+
 \frac\theta N[\vect n]_\times\right)^N,
\]
并得到罗德里格公式。
\end{advexamplebox}
\begin{advsolution}
因 $[\vect n]_\times^2=\vect n\vect n^{\mathsf T}-\mathsf I$，且
$[\vect n]_\times^3=-[\vect n]_\times$，指数级数的奇、偶次项可分别求和：
\[
 e^{\theta[\vect n]_\times}
 =\mathsf I+\sin\theta[\vect n]_\times
 +(1-\cos\theta)[\vect n]_\times^2.
\]
作用于任意 $\vect r$ 便得
\[
 R_{\vect n}(\theta)\vect r
 =\vect r\cos\theta+(\vect n\times\vect r)\sin\theta
 +\vect n(\vect n\cdot\vect r)(1-\cos\theta),
\]
即罗德里格公式。极限形式是矩阵指数的标准定义，表明有限转动可看成无穷多个无穷小转动的有序乘积。
\end{advsolution}

\begin{advexamplebox}{角速度矩阵与伪矢量的同变换}\label{advA1:ex:omega-transform}
对 $\mathsf O\in SO(3)$，证明
\[
 [\vect\omega']_\times=\mathsf O[\vect\omega]_\times\mathsf O^{-1}
 \quad\Longleftrightarrow\quad
 \vect\omega'=\mathsf O\vect\omega .
\]
\end{advexamplebox}
\begin{advsolution}
对任意矢量 $\vect x$，
\[
 \mathsf O[\vect\omega]_\times\mathsf O^{-1}\vect x
 =\mathsf O(\vect\omega\times\mathsf O^{-1}\vect x)
 =(\mathsf O\vect\omega)\times\vect x.
\]
第二个等号用了 $\det\mathsf O=+1$ 的保向性。因为等式对所有 $\vect x$ 成立，所以两个叉乘矩阵相等。反向推理同理。若把 $SO(3)$ 放宽为包含镜反射的 $O(3)$，就会多出 $\det\mathsf O$，这正是角速度称为伪矢量的原因。
\end{advsolution}

\subsubsection{欧拉角与角速度}

采用 $Z-X_1-Z_2$ 约定：先绕定系 $Z$ 轴转进动角 $\varphi$，再绕节线 $X_1$ 转章动角 $\theta$，最后绕刚体轴 $Z_2=z$ 转自转角 $\psi$。
\begin{center}\begin{minipage}{\linewidth}\centering
\begin{tikzpicture}[scale=.87]
 \begin{scope}[shift={(0,0)}]
  \draw[axis] (0,0)--(1.35,0) node[right]{$X$};
  \draw[curve,->] (0,0)--(1.10,.64) node[right]{$X_1$};
  \draw[curve,->] (0,0)--(0,1.35) node[above]{$Z_1=Z$};
  \draw[->,TextGray] (.48,.05) arc[start angle=6,end angle=30,radius=.50];
  \node at (.25,-1.05) {\strut(1)};
 \end{scope}
 \begin{scope}[shift={(5,0)}]
  \draw[curve,->] (0,0)--(1.35,0) node[below]{$X_2=X_1$};
  \draw[axis] (0,0)--(0,1.35) node[above]{$Z_1$};
  \draw[curve,->] (0,0)--(.67,1.16) node[above right]{$Z_2$};
  \draw[->,TextGray] (.05,.50) arc[start angle=84,end angle=60,radius=.50];
  \node at (.25,-1.05) {\strut(2)};
 \end{scope}
 \begin{scope}[shift={(10,0)}]
  \draw[axis] (0,0)--(1.35,0) node[right]{$X_2$};
  \draw[curve,->] (0,0)--(1.10,.64) node[right]{$x$};
  \draw[curve,->] (0,0)--(.30,1.32) node[above]{$z=Z_2$};
  \draw[->,TextGray] (.48,.05) arc[start angle=6,end angle=30,radius=.50];
  \node at (.25,-1.05) {\strut(3)};
 \end{scope}
 \draw[-{Latex[length=2.2mm]},TextGray,thick] (2.05,.20)--(3.60,.20);
 \draw[-{Latex[length=2.2mm]},TextGray,thick] (7.05,.20)--(8.60,.20);
\end{tikzpicture}
\bookcaption{欧拉角的三步有序转动；次序改变会改变最终姿态}\label{advA1:fig:euler-steps}
\end{minipage}\end{center}
图\ref{advA1:fig:euler-steps}将三次转动分开画出，避免在一张三维图中混淆中间轴。由矩阵乘积的莱布尼茨求导法则，总角速度是三步瞬时角速度的矢量和。在刚体主轴系中，
\begin{equation}\label{advA1:eq:euler-kinematic}
 \begin{aligned}
 \omega_1&=\dot\varphi\sin\theta\sin\psi+\dot\theta\cos\psi,\\
 \omega_2&=\dot\varphi\sin\theta\cos\psi-\dot\theta\sin\psi,\\
 \omega_3&=\dot\varphi\cos\theta+\dot\psi.
 \end{aligned}
\end{equation}
这是欧拉运动学方程。在 $\theta=0$ 处 $\varphi$ 与 $\psi$ 的作用重合，是坐标奇点，不是刚体姿态的物理奇点。

\subsection{惯量张量与欧拉动力学方程}

对绕定点以角速度 $\vect\omega$ 转动的刚体，
\[
 T=\frac12\int\rho(\vect\omega\times\vect r)^2\dd V,
 \qquad
 \vect L=\int\rho\vect r\times(\vect\omega\times\vect r)\dd V.
\]
用三重矢积公式，
\begin{equation}\label{advA1:eq:inertia-tensor}
 \vect L=\mathsf I\vect\omega,
 \quad
 \mathsf I=\int\rho(r^2\mathsf 1-\vect r\vect r^{\mathsf T})\dd V,
 \quad
 T=\frac12\vect\omega\cdot\mathsf I\vect\omega.
\end{equation}
若二阶质量矩矩阵为 $\mathsf M=\int\rho\vect r\vect r^{\mathsf T}\dd V$，则
$\mathsf I=\operatorname{tr}(\mathsf M)\mathsf1-\mathsf M$。

\begin{advexamplebox}{绕任意轴的转动惯量}\label{advA1:ex:axis-inertia}
刚体绕经定点的单位方向 $\vect n$ 以角速度 $\omega$ 转动。证明该轴的转动惯量为
\[
 I_{\vect n}=\vect n^{\mathsf T}\mathsf I\vect n.
\]
\end{advexamplebox}
\begin{advsolution}
令 $\vect\omega=\omega\vect n$，将式\eqref{advA1:eq:inertia-tensor}中的动能与定义 $T=I_{\vect n}\omega^2/2$ 比较，
\[
 \frac12I_{\vect n}\omega^2
 =\frac12\omega^2\vect n^{\mathsf T}\mathsf I\vect n.
\]
因而得证。同时
$\vect L\cdot\vect n=\omega\vect n^{\mathsf T}\mathsf I\vect n=I_{\vect n}\omega$，但若 $\vect n$ 不是主轴，$\vect L$ 本身不与 $\vect\omega$ 平行。
\end{advsolution}

\begin{advtheorem}{主轴定理}\label{advA1:thm:principal-axes}
实对称惯量张量存在一组正交主轴，在该基底下
\[
 \mathsf I=\operatorname{Diag}(I_1,I_2,I_3).
\]
\end{advtheorem}
这是实对称矩阵的谱定理。解
$\det(\mathsf I-\lambda\mathsf1)=0$得三个实本征值；不同本征值的本征矢量自动正交，简并子空间内再作格拉姆--施密特正交化。以三个归一化本征矢量为列构造正交矩阵，即得对角化。在主轴系中
\[
 L_i=I_i\omega_i,
 \qquad T=\frac12(I_1\omega_1^2+I_2\omega_2^2+I_3\omega_3^2).
\]

\begin{advtheorem}{欧拉动力学方程}\label{advA1:thm:euler-dynamics}
在刚体主轴系中，外力矩分量为 $M_i$ 时，
\begin{equation}\label{advA1:eq:euler-dynamics}
 \begin{aligned}
 I_1\dot\omega_1+(I_3-I_2)\omega_2\omega_3&=M_1,\\
 I_2\dot\omega_2+(I_1-I_3)\omega_3\omega_1&=M_2,\\
 I_3\dot\omega_3+(I_2-I_1)\omega_1\omega_2&=M_3.
 \end{aligned}
\end{equation}
\end{advtheorem}
证明：惯性系中 $\dd\vect L/\dd t=\vect M$。任意矢量的两系导数关系为
\[
 \left(\frac{\dd\vect L}{\dd t}\right)_{\rm inertial}
 =\left(\frac{\dd\vect L}{\dd t}\right)_{\rm body}
 +\vect\omega\times\vect L.
\]
在主轴系中 $\vect L=(I_1\omega_1,I_2\omega_2,I_3\omega_3)$，逐分量展开即得式\eqref{advA1:eq:euler-dynamics}。

把式\eqref{advA1:eq:euler-kinematic}代入动能，就可用欧拉角写拉格朗日方程。外力矩的功率为
\[
 \vect M\cdot\vect\omega=M_\varphi\dot\varphi
 +M_\theta\dot\theta+M_\psi\dot\psi,
\]
其中
\[
 M_\varphi=M_1\sin\theta\sin\psi+M_2\sin\theta\cos\psi+M_3\cos\theta,
\]
\[
 M_\theta=M_1\cos\psi-M_2\sin\psi,
 \qquad M_\psi=M_3.
\]
它们就是三个欧拉角对应的广义力。

\subsubsection{自由非对称刚体}

设 $I_1<I_2<I_3$。无外力矩时，角动量大小与动能守恒。记刚体系中的角动量分量为
$\Pi_i=I_i\omega_i$，则轨道是
\begin{equation}\label{advA1:eq:pi-intersection}
 \Pi_1^2+\Pi_2^2+\Pi_3^2=L^2,
 \qquad
 \frac{\Pi_1^2}{2I_1}+\frac{\Pi_2^2}{2I_2}
 +\frac{\Pi_3^2}{2I_3}=T.
\end{equation}
即角动量球与能量椭球的交线。
\begin{center}\begin{minipage}{\linewidth}\centering
\begin{tikzpicture}[scale=.88]
 \begin{scope}
  \draw[axis] (-2.3,0)--(2.4,0) node[right]{$\Pi_1$};
  \draw[axis] (0,-2.0)--(0,2.1) node[above]{$\Pi_3$};
  \foreach \sgn in {-1,1}{
   \draw[curve,domain=0:360,samples=121,variable=\t] plot ({.83666*cos(\t)+.460163*sin(\t)},{\sgn*1.8*sqrt(.3+.525*cos(\t)^2)+.167332*sin(\t)});
   \fill (0,{\sgn*1.8}) circle(1.8pt);
   \fill ({\sgn*2},0) circle(1.8pt);
  }
  \node at (0,-2.42){(a) 交线投影；点为稳定主轴方向};
 \end{scope}
 \begin{scope}[shift={(6,0)}]
  \draw[very thick] (0,0) ellipse (1.55 and 1.02);
  \draw[ground] (-2,-1.02)--(2,-1.02);
  \draw[vector] (0,0)--(0,-1.75) node[right]{$\vect L$};
  \fill (0,-1.02) circle(2pt);
  \draw[thin] (.10,-1.02)--(.65,-.70) node[right]{接触点};
  \fill (0,0) circle(1.4pt) node[above]{$O$};
  \node at (0,1.45){能量椭球};
  \node at (0,-2.42){(b) 定平面垂直于 $\vect L$};
 \end{scope}
\end{tikzpicture}
\bookcaption{自由刚体的角动量交线与能量椭球纯滚动图像}\label{advA1:fig:poinsot}
\end{minipage}\end{center}
图\ref{advA1:fig:poinsot}(a) 是三维交线的斜投影，未画出 $\Pi_2$ 轴；投影曲线的交叉不表示三维轨道相交。绕最小和最大主轴的转动稳定，绕中间主轴的转动不稳定，后者会表现为反复翻转，即所谓中间轴定理或扎尼别科夫效应。在角速度空间，能量椭球在惯性系中沿一个与守恒 $\vect L$ 垂直的定平面纯滚动，如图\ref{advA1:fig:poinsot}(b)，这就是庞锁构造。

\subsubsection{对称陀螺与章动}

对 $I_1=I_2=I, I_3=I_0$ 的对称陀螺，
\begin{equation}\label{advA1:eq:symmetric-top-T}
 T=\frac I2(\dot\theta^2+\dot\varphi^2\sin^2\theta)
 +\frac{I_0}{2}(\dot\psi+\dot\varphi\cos\theta)^2.
\end{equation}
若势能仅为 $V(\theta)$，$\varphi,\psi$ 都是循环坐标：
\[
 J_\varphi=I\dot\varphi\sin^2\theta+J_\psi\cos\theta,
 \qquad
 J_\psi=I_0(\dot\psi+\dot\varphi\cos\theta).
\]
同时罗斯约化两个循环坐标，
\begin{equation}\label{advA1:eq:top-effective}
 \mathcal R=\frac I2\dot\theta^2-V_{\rm eff}(\theta)-\frac{J_\psi^2}{2I_0},
 \quad
 V_{\rm eff}=\frac{(J_\varphi-J_\psi\cos\theta)^2}
 {2I\sin^2\theta}+V(\theta).
\end{equation}
定义约化能 $E_\theta=E-J_\psi^2/(2I_0)$，转向点由 $E_\theta=V_{\rm eff}$ 决定，再代回
\[
 \dot\varphi=\frac{J_\varphi-J_\psi\cos\theta}{I\sin^2\theta},
 \qquad
 \dot\psi=\frac{J_\psi}{I_0}-\dot\varphi\cos\theta
\]
即可求进动与自转。
\begin{center}\begin{minipage}{\linewidth}\centering
\begin{tikzpicture}[x=1.9cm,y=1.45cm]
 \draw[axis] (0,0)--(3.35,0) node[right]{$\theta$};
 \draw[axis] (0,0)--(0,2.35) node[above]{$V_{\rm eff}$};
 \draw[curve,domain=.50:2.64159,samples=101] plot (\x,{.5/(sin(\x r)^2)});
 \draw[TextGray,thick] (.4,1)--(2.75,1) node[right]{$E_\theta$};
 \foreach \x/\lab in {.785398/\theta_-,2.356194/\theta_+}{\draw[guide] (\x,0)--(\x,1);\fill (\x,1) circle(1.6pt);\node[below] at (\x,0){$\lab$};}
 \fill (1.570796,.5) circle(1.7pt) node[below=4pt]{$\theta_0$};
 \node[below] at (3.14159,0){$\pi$};
 \node[font=\small] at (1.6,2.7){示例：$V_{\rm eff}=1/(2\sin^2\theta)$};
\end{tikzpicture}
\bookcaption{对称陀螺的有效势、稳定无章动点与章动转向点}\label{advA1:fig:top-effective}
\end{minipage}\end{center}
图\ref{advA1:fig:top-effective}中 $\theta_0$ 是有效势极小点；能量稍高时，$\theta$ 在 $\theta_-$ 与 $\theta_+$ 之间章动。

\begin{advexamplebox}{无章动条件与小章动频率}\label{advA1:ex:nutation}
对式\eqref{advA1:eq:top-effective}的一般对称陀螺，求恒定 $\theta=\theta_0$ 的条件，并求其附近小章动的频率。
\end{advexamplebox}
\begin{advsolution}
无章动要求
\[
 V_{\rm eff}'(\theta_0)=0,
 \qquad E=V_{\rm eff}(\theta_0)+\frac{J_\psi^2}{2I_0}.
\]
令 $A(\theta)=J_\varphi-J_\psi\cos\theta$，第一个条件显式写成
\[
 \frac1I\left[
 \frac{A J_\psi}{\sin\theta}
 -\frac{A^2\cos\theta}{\sin^3\theta}\right]_{\theta_0}
 +V'(\theta_0)=0.
\]
令 $\theta=\theta_0+\eta$，将有效势展开到二阶：
\[
 I\ddot\eta+V_{\rm eff}''(\theta_0)\eta=0.
\]
因而
\[
 \boxed{\omega_{\rm nut}^2=\frac{V_{\rm eff}''(\theta_0)}I}.
\]
稳定无章动还要求 $V_{\rm eff}''(\theta_0)>0$；若小于零，则微小扰动指数增长。
\end{advsolution}

\begin{advexamplebox}{始终保持底边水平的等边三角板}\label{advA1:ex:triangle-top}
均匀等边三角形薄板 $OAB$ 的质量为 $m$、边长为 $l$，绕过顶点 $O$ 的竖直轴定点转动，且 $AB$ 始终水平。
(1) 当板以进动角速度 $\omega$ 作无章动运动时，求板中线 $OG$ 与竖直向下方向的夹角 $\phi$；
(2) 在不施加额外竖直轴力矩的情况下，求 $\phi$ 发生微小章动的周期。
\end{advexamplebox}
\AdvHideHeadingtrue
\begin{advsolution}
\noindent\begin{minipage}{\linewidth}{\bfseries\color{DeepBlue}解答}\par
\begin{center}\begin{minipage}{\linewidth}\centering
\begin{tikzpicture}[scale=.92]
 \draw[axis] (0,1.9)--(0,-2.0) node[below]{竖直轴};
 \coordinate (O) at (0,.85);
 \coordinate (G) at (.7,-.083333);
 \coordinate (A) at (.15,-.55);
 \coordinate (B) at (1.95,-.55);
 \fill[DeepBlue!7] (O)--(A)--(B)--cycle;
 \draw[very thick] (O)--(A)--(B)--cycle;
 \draw[guide] (O)--(0,-1.35);
 \draw[curve] (O)--(G) node[midway,above right]{$OG$};
 \fill (G) circle(2.3pt) node[right=4pt]{$G$};
 \draw[->,DeepBlue] ($(O)+(0,-.48)$) arc[start angle=-90,end angle=-53,radius=.48];
 \node at (.40,.20) {$\phi$};
 \draw[very thick,TextGray] (A)--(B) node[midway,below=4pt]{$AB$水平};
 \draw[->,TextGray] (-.55,.85) arc[start angle=180,end angle=20,x radius=.85,y radius=.25];
 \node at (-.82,1.15) {$\alpha$};
\end{tikzpicture}
\bookcaption{等边三角板的中线、水平底边和进动角}\label{advA1:fig:triangle}
\end{minipage}\end{center}\end{minipage}

如图\ref{advA1:fig:triangle}，记板的方位角为 $\alpha$。质心到 $O$ 的距离
\[
 a=OG=\frac{l}{\sqrt3}.
\]
取过 $O$ 且平行 $AB$ 的轴、板法线和 $OG$ 为三条主轴，其主惯量为
\[
 I_1=\frac{3}{8}ml^2,
 \qquad I_2=\frac{5}{12}ml^2,
 \qquad I_3=\frac1{24}ml^2,
 \qquad I_2-I_3=I_1.
\]
动能和势能为
\[
 T=\frac12I_1\dot\phi^2+
 \frac12A(\phi)\dot\alpha^2,
 \quad A(\phi)=I_2\sin^2\phi+I_3\cos^2\phi,
 \qquad V=-mga\cos\phi.
\]

(1) 无章动时 $\dot\alpha=\omega$，$\phi$ 方程给出
\[
 (I_2-I_3)\omega^2\sin\phi\cos\phi-mga\sin\phi=0.
\]
非竖直解为
\[
 \boxed{\cos\phi=\frac{mga}{(I_2-I_3)\omega^2}
 =\frac{8g}{3\sqrt3\,l\omega^2}}.
\]
它要求 $\omega^2\ge8g/(3\sqrt3\,l)$。

(2) 章动时 $\alpha$ 不再保持严格匀速，守恒的是
$p_\alpha=A(\phi)\dot\alpha=A(\phi_0)\omega$。罗斯函数为
\[
 \mathcal R=\frac12I_1\dot\phi^2+mga\cos\phi
 -\frac{p_\alpha^2}{2A(\phi)}.
\]
在平衡角 $\phi_0=\phi$ 附近线性化，并用第(1)问的平衡条件化简，得
\[
 \Omega_{\rm nut}^2
 =\omega^2\sin^2\phi
 \left[1+\frac{4I_1\cos^2\phi}{A(\phi)}\right]
 =\omega^2\sin^2\phi\,
 \frac{10+27\cos^2\phi}{10-9\cos^2\phi}.
\]
因而
\[
 \boxed{T_{\rm nut}=\frac{2\pi}{\omega\sin\phi}
 \sqrt{\frac{10-9\cos^2\phi}{10+27\cos^2\phi}}}.
\]
若误把章动过程中的 $\dot\alpha$ 仍当作外部强制的常量，会遗漏根号中的惯性耦合因子。
\end{advsolution}

\section{连续体}

\subsection{分析力学的变分表述}

有限维函数的极值由导数为零刻画；当未知对象本身是函数时，研究对象就变成泛函。设区域
$D\subset\mathbb R^n$ 上有若干函数 $f_\mu(\vect x)$，局域且只含一阶导数的密度泛函写成
\[
 \mathcal L=\mathcal L(f_\mu,\partial_i f_\mu,x_i),
 \qquad
 S[f_\mu]=\int_D\mathcal L\,\dd^n x .
\]

\begin{advdefinition}{泛函导数}\label{advA1:def:functional-derivative}
若对满足指定边界条件的任意测试函数 $g_\mu$，
\[
 \left.\frac{\dd}{\dd\epsilon}S[f_\mu+\epsilon g_\mu]\right|_{\epsilon=0}
 =\int_D\sum_\mu\frac{\delta S}{\delta f_\mu}g_\mu\,\dd^n x ,
\]
则称 $\delta S/\delta f_\mu$ 为泛函在 $f_\mu$ 处的泛函导数。
\end{advdefinition}

\begin{advtheorem}{欧拉--拉格朗日方程}\label{advA1:thm:field-el}
当边界值 $f_\mu|_{\partial D}$ 固定时，$S$ 的一阶变分为零等价于
\begin{equation}\label{advA1:eq:field-el}
 \frac{\delta S}{\delta f_\mu}
 =\frac{\partial\mathcal L}{\partial f_\mu}
 -\sum_i\partial_i\left(
 \frac{\partial\mathcal L}{\partial(\partial_i f_\mu)}\right)=0 .
\end{equation}
\end{advtheorem}
证明：令 $\delta f_\mu=\epsilon g_\mu$，展开到一阶：
\[
 \delta S=\int_D\sum_\mu\left[
 \frac{\partial\mathcal L}{\partial f_\mu}g_\mu+
 \sum_i\frac{\partial\mathcal L}{\partial(\partial_i f_\mu)}
 \partial_i g_\mu\right]\dd^n x .
\]
对第二项分部积分。边界上 $g_\mu=0$，边界项消失；区域内 $g_\mu$ 任意，故其系数必须处处为零，即得式\eqref{advA1:eq:field-el}。

这个推导同时给出三个常用变式。若边界值自由，分部积分留下自然边界条件
\[
 \left.\vect n\cdot
 \frac{\partial\mathcal L}{\partial(\nabla f_\mu)}\right|_{\partial D}=0.
\]
若密度含到 $n$ 阶导数，则反复分部积分得到欧拉--泊松方程
\[
 \sum_{k=0}^{n}(-1)^k\partial_{i_1}\cdots\partial_{i_k}
 \left[\frac{\partial\mathcal L}
 {\partial(\partial_{i_1}\cdots\partial_{i_k}f_\mu)}\right]=0.
\]
若还带有积分约束 $C_a[f]=0$，则对
$\mathcal L-\sum_a\lambda_a\mathcal C_a$ 使用同一方程，并与全部约束联立求 $\lambda_a$。

\begin{advexamplebox}{固定均值与能量时的最大熵分布}\label{advA1:ex:max-entropy}
求一维速度分布 $f(v)$，使
\[
 -\int_{\mathbb R}f\ln f\,\dd v
\]
在
\[
 \int f\,\dd v=1,\qquad \int vf\,\dd v=0,\qquad
 \int\frac12mv^2f\,\dd v=\frac12kT
\]
条件下最大。
\end{advexamplebox}
\begin{advsolution}
引入三个拉格朗日乘数，令
\[
 \mathcal F[f]=-\int f\ln f\,\dd v-\lambda_0\!\left(\int f\,\dd v-1\right)
 -\lambda_1\int vf\,\dd v
 -\lambda_2\!\left(\int\frac12mv^2f\,\dd v-\frac12kT\right).
\]
泛函导数为
\[
 \frac{\delta\mathcal F}{\delta f}
 =-(\ln f+1)-\lambda_0-\lambda_1v-\frac12\lambda_2mv^2=0.
\]
因此 $f=Ae^{-av^2-bv}$。零均值给出 $b=0$，归一化与能量约束给出
\[
 \boxed{f(v)=\sqrt{\frac{m}{2\pi kT}}
 \exp\!\left(-\frac{mv^2}{2kT}\right)}.
\]
二阶变分
$\delta^2(-\int f\ln f)=-\int(\delta f)^2/f\,\dd v<0$，故驻点确为最大值。其微分熵为
\[
 S_{\max}=\frac12\left[1+\ln\!\left(\frac{2\pi kT}{m}\right)\right].
\]
\end{advsolution}

\newcommand{\AOneElasticSecond}{%
\subsubsection{应力、各向同性本构关系与弹性波}

过一点取有向面元 $\dd\vect S=\vect n\,\dd S$。相邻两侧材料通过该面元传递的力在局域极限下与面元线性：
\[
 \dd\vect F=\bm\sigma\,\dd\vect S,
\]
其中 $\bm\sigma$ 是应力张量。对角元是正应力，非对角元是剪应力。

\begin{advtheorem}{经典连续体的应力张量对称}\label{advA1:thm:stress-symmetric}
若质元内部不存在独立的偶应力，局域角动量守恒要求
\[
 \boxed{\sigma_{ij}=\sigma_{ji}}.
\]
\end{advtheorem}
证明：取边长为 $\ell$ 的微小立方体。相对两个面上的应力合力对质心的力矩中，体力矩是
$\order{\ell^4}$，而剪应力对产生的力偶是 $\order{\ell^3}$。例如绕 $z$ 轴的领先项为
$(\sigma_{21}-\sigma_{12})\ell^3$。令 $\ell\to0$ 时角动量平衡必须逐阶成立，故
$\sigma_{12}=\sigma_{21}$；循环指标即得其余分量。
\begin{center}\begin{minipage}{\linewidth}\centering
\begin{tikzpicture}[scale=.92]
 \draw[very thick] (0,0) rectangle (2.2,1.6);
 \draw[vector] (0,1.25)--(0,.45) node[midway,left]{$\sigma_{21}$};
 \draw[vector] (2.2,.35)--(2.2,1.15) node[midway,right]{$\sigma_{21}$};
 \draw[vector] (1.25,0)--(.45,0) node[midway,below]{$\sigma_{12}$};
 \draw[vector] (.95,1.6)--(1.75,1.6) node[midway,above]{$\sigma_{12}$};
 \node at (1.1,.8) {微小质元};
\end{tikzpicture}
\bookcaption{剪应力不成对相等时会留下随体积缩放的净力矩}\label{advA1:fig:stress-symmetry}
\end{minipage}\end{center}
图\ref{advA1:fig:stress-symmetry}中的箭头作用在不同面上；相等的是互易剪应力，而不是把同一面上的两个方向混为一谈。

\begin{advtheorem}{圣维南--基尔霍夫各向同性线弹性关系}\label{advA1:thm:svk}
各向同性、线性小应变材料的弹性能密度和应力为
\begin{equation}\label{advA1:eq:svk}
 W=\frac{\lambda}{2}(\operatorname{tr}\bm\varepsilon)^2
 +\mu\,\operatorname{tr}(\bm\varepsilon^2),
 \qquad
 \bm\sigma=\lambda(\operatorname{tr}\bm\varepsilon)\mathsf1
 +2\mu\bm\varepsilon ,
\end{equation}
其中 $\lambda,\mu$ 是拉梅参数。
\end{advtheorem}
证明：二阶对称张量在 $SO(3)$ 下分为迹和对称无迹两个互不混合的不可约部分。各向同性二次型只能分别给这两个部分赋予一个标量系数，不能出现依赖方向的交叉项，故
$W=a(\operatorname{tr}\bm\varepsilon)^2+b\,\operatorname{tr}(\bm\varepsilon^2)$。把常数重命名为
$a=\lambda/2,b=\mu$，再用
$\sigma_{ij}=\partial W/\partial\varepsilon_{ij}$ 即得式\eqref{advA1:eq:svk}。

线弹性场的拉格朗日密度是
\[
 \mathcal L=\frac12\rho_m\dot{\vect\xi}^{\,2}-W,
\]
其中 $\rho_m$ 是质量密度。场方程为纳维--柯西方程
\begin{equation}\label{advA1:eq:navier-cauchy}
 \rho_m\ddot{\vect\xi}
 =(\lambda+\mu)\nabla(\nabla\cdot\vect\xi)
 +\mu\nabla^2\vect\xi+\vect f .
\end{equation}

\begin{advexamplebox}{由应力散度推出纳维--柯西方程}\label{advA1:ex:stress-divergence}
从应力张量的局域动量平衡
$\rho_m\ddot{\vect\xi}=\nabla\cdot\bm\sigma+\vect f$
出发，推导式\eqref{advA1:eq:navier-cauchy}。
\end{advexamplebox}
\begin{advsolution}
由式\eqref{advA1:eq:svk}及
$2\bm\varepsilon=\nabla\vect\xi+(\nabla\vect\xi)^{\mathsf T}$，
\[
 \bm\sigma=\lambda(\nabla\cdot\vect\xi)\mathsf1
 +\mu[\nabla\vect\xi+(\nabla\vect\xi)^{\mathsf T}].
\]
逐分量取散度：
\[
 (\nabla\cdot\bm\sigma)_i
 =\lambda\partial_i\partial_j\xi_j
 +\mu\partial_j\partial_j\xi_i
 +\mu\partial_j\partial_i\xi_j.
\]
交换光滑场的偏导次序，得
\[
 \nabla\cdot\bm\sigma
 =(\lambda+\mu)\nabla(\nabla\cdot\vect\xi)
 +\mu\nabla^2\vect\xi,
\]
代回动量平衡即得所求。边界牵引条件是
$\bm\sigma\cdot\vect n=\dd\vect F/\dd S$。
\end{advsolution}

令 $\vect\xi=\vect A e^{\ii(\vect k\cdot\vect r-\omega t)}$，把振幅分解为纵向
$\vect A_\parallel$ 与横向 $\vect A_\perp$。式\eqref{advA1:eq:navier-cauchy}给出
\[
 c_{\rm L}=\sqrt{\frac{\lambda+2\mu}{\rho_m}},
 \qquad
 c_{\rm T}=\sqrt{\frac{\mu}{\rho_m}}.
\]
纵波只有一个偏振，横波有两个彼此正交的偏振；因此固体中的首波（P 波）比次波（S 波）快。
\begin{center}\begin{minipage}{\linewidth}\centering
\begin{tikzpicture}[scale=.9]
 \begin{scope}
  \draw[axis] (-.2,0)--(4.3,0) node[right]{$\vect k$};
  \foreach \x in {0,.46,1.02,1.72,2.48,3.12,3.60}{\fill (\x,0) circle(2pt);}
  \draw[<->,DeepBlue,thick] (1.2,-.38)--(1.75,-.38);
  \node[DeepBlue,left] at (1.0,-.38) {$\vect A_\parallel$};
  \node at (2,-.9) {纵波};
 \end{scope}
 \begin{scope}[shift={(6.0,0)}]
  \draw[axis] (-.2,0)--(4.3,0) node[right]{$\vect k$};
  \foreach \x in {0,.6,...,3.6}{\fill (\x,{.45*sin(120*\x)}) circle(2pt);}
  \draw[<->,DeepBlue,thick] (1.2,-.35)--(1.2,.45) node[right]{$\vect A_\perp$};
  \node at (2,-.9) {横波};
 \end{scope}
\end{tikzpicture}
\bookcaption{弹性体纵波的位移平行于波矢，横波的位移垂直于波矢}\label{advA1:fig:elastic-waves}
\end{minipage}\end{center}
图\ref{advA1:fig:elastic-waves}把传播方向与位移方向同时标出，二者不能仅凭波形外观判断。

单轴拉伸定义杨氏模量 $E_Y$ 与泊松比 $\nu$。它们和拉梅参数的关系为
\[
 \nu=\frac{\lambda}{2(\lambda+\mu)},\qquad
 E_Y=\frac{\mu(3\lambda+2\mu)}{\lambda+\mu},
\]
\[
 \lambda=\frac{\nu E_Y}{(1+\nu)(1-2\nu)},\qquad
 \mu=\frac{E_Y}{2(1+\nu)}.
\]
体积模量与剪切模量为
\[
 K=\lambda+\frac{2\mu}{3},\qquad G=\mu.
\]
其中
$\operatorname{tr}\bm\sigma=3K\operatorname{tr}\bm\varepsilon$，而偏应力与偏应变满足
$\bm\Sigma=2G\bm E$。

\begin{advexamplebox}{梁的均匀弯曲}\label{advA1:ex:beam-bending}
长为 $L$、宽为 $b$、厚为 $h$ 的均匀矩形梁弯成总转角 $\Theta$ 的圆弧。质量密度为
$\rho_m$。求应变分布、弹性能和维持弯曲所需端力矩，并用拉梅参数表示结果。
\end{advexamplebox}
\begin{advsolution}
欧拉--伯努利假设下，中性面到距离 $y$ 处的纵向应变为
\[
 \varepsilon_{xx}=-\kappa y,\qquad \kappa=\frac{\Theta}{L}.
\]
侧面自由时应力是单轴的，
$\sigma_{xx}=E_Y\varepsilon_{xx}$。截面二次矩
$I_b=bh^3/12$，所以
\[
 M=\int_A(-y)\sigma_{xx}\dd A
 =E_YI_b\kappa
 =\boxed{\frac{E_Ybh^3}{12L}\Theta}.
\]
弹性能
\[
 U=\frac12\int_0^L E_YI_b\kappa^2\dd x
 =\boxed{\frac{E_Ybh^3}{24L}\Theta^2}.
\]
若用质量 $m=\rho_mbhL$ 消去 $b$，
\[
 U=\frac{E_Ymh^2}{24\rho_mL^2}\Theta^2,
 \qquad
 E_Y=\frac{\mu(3\lambda+2\mu)}{\lambda+\mu}.
\]
厚度以三次方进入弯曲刚度，这是最容易在尺度估算中遗漏的关系。
\begin{center}\begin{minipage}{\linewidth}\centering
\begin{tikzpicture}[scale=.9]
 \coordinate (C) at (0,0);
 \fill[DeepBlue!8] (35:2.7) arc(35:145:2.7)--(145:3.3) arc(145:35:3.3)--cycle;
 \draw[very thick] (35:2.7) arc(35:145:2.7)--(145:3.3) arc(145:35:3.3)--cycle;
 \draw[guide] (35:3) arc(35:145:3);
 \draw[guide] (C)--(35:3.6) (C)--(145:3.6);
 \draw[DeepBlue,->] (35:.85) arc(35:145:.85);
 \node at (0,1.08){$\Theta$};\fill(C) circle(1.5pt) node[below]{$C$};
 \draw[<->] (C)--(90:3) node[midway,right]{$R$};
 \draw[<->] (35:2.7)--(35:3.3) node[midway,right=5pt]{$h$};
 \node at (0,3.75){同心弯梁截面};
\end{tikzpicture}
\bookcaption{梁弯曲时中性层不伸长，上下表面的纵向应变符号相反}\label{advA1:fig:beam}
\end{minipage}\end{center}
\end{advsolution}

\begin{advexamplebox}{圆柱的扭转}\label{advA1:ex:cylinder-torsion}
半径 $R$、长度 $L$ 的均匀实心圆柱，底面固定，顶面刚性转过小角 $\Theta$。求位移场、剪应力、维持转动所需力矩和弹性能。
\end{advexamplebox}
\begin{advsolution}
取柱坐标，满足两端边界的位移为
\[
 \xi_\varphi(r,z)=\frac{\Theta z}{L}r,\qquad
 \xi_r=\xi_z=0.
\]
工程剪应变与剪应力为
\[
 \gamma_{z\varphi}=\frac{\partial\xi_\varphi}{\partial z}
 =\frac{\Theta r}{L},\qquad
 \tau_{z\varphi}=\mu\frac{\Theta r}{L}.
\]
顶面力矩
\[
 M=\int_A\tau_{z\varphi}r\,\dd A
 =\frac{\mu\Theta}{L}\int_0^Rr^2(2\pi r\dd r)
 =\boxed{\frac{\pi\mu R^4}{2L}\Theta}.
\]
因此
\[
 \boxed{U=\frac12M\Theta
 =\frac{\pi\mu R^4}{4L}\Theta^2}.
\]
空心圆柱只需把 $R^4$ 换为 $R_{\rm o}^4-R_{\rm i}^4$。力矩量纲
$[\mu R^4/L]=FL$，检查正确。
\begin{center}\begin{minipage}{\linewidth}\centering
\begin{tikzpicture}[scale=.86]
 \draw[very thick] (-1.2,-1.5)--(-1.2,1.35);
 \draw[very thick] (1.2,-1.5)--(1.2,1.35);
 \draw[very thick] (0,1.35) ellipse (1.2 and .35);
 \draw[very thick] (0,-1.5) ellipse (1.2 and .35);
 \draw[->,DeepBlue,thick] (.20,1.85) arc[start angle=105,end angle=-25,x radius=1.15,y radius=.42];
 \node at (1.45,1.72) {$\Theta$};
 \draw[<->,TextGray] (1.55,-1.5)--(1.55,1.35) node[midway,right]{$L$};
 \draw[ground] (-1.55,-1.86)--(1.55,-1.86);
\end{tikzpicture}
\bookcaption{圆柱顶面相对底面的刚性扭转与轴向长度}\label{advA1:fig:torsion}
\end{minipage}\end{center}
\end{advsolution}

}

\newcommand{\AOneFluidSection}{%
\subsection{流体}

流体没有由静态形变直接决定的恢复应力，其应力取决于形变速率，并且常与热学自由度耦合。因此用欧拉描述：速度、密度、压强和温度是空间与时间的场
$\vect v(\vect r,t),\rho(\vect r,t),p(\vect r,t),T(\vect r,t)$。质量守恒为
\begin{equation}\label{advA1:eq:continuity}
 \frac{\partial\rho}{\partial t}+\nabla\cdot(\rho\vect v)=0.
\end{equation}

\begin{advdefinition}{随体全导数}\label{advA1:def:material-derivative}
\[
 \frac{D}{Dt}=\frac{\partial}{\partial t}+\vect v\cdot\nabla.
\]
它是在跟随同一流体质元运动时测得的变化率。
\end{advdefinition}
由式\eqref{advA1:eq:continuity}，
\[
 \frac{D\rho}{Dt}=-\rho\nabla\cdot\vect v.
\]
速度梯度分解为
\[
 \nabla\vect v=\bm U+\bm W,\qquad
 \bm U=\frac12[\nabla\vect v+(\nabla\vect v)^{\mathsf T}],
 \quad
 \bm W=\frac12[\nabla\vect v-(\nabla\vect v)^{\mathsf T}],
\]
其中 $\bm U$ 是应变率张量，$\bm W$ 是涡张量；涡量矢量为
$\vect\omega=\nabla\times\vect v$。再去掉 $\bm U$ 的迹得到偏应变率
$\bm U'=\bm U-(\nabla\cdot\vect v)\mathsf1/3$。

\begin{advdefinition}{线性本构前提}\label{advA1:def:linear-constitutive}
本构方程中的材料参数视为与流体自由度无关的常数；应力对形变率只保留线性项。
\end{advdefinition}

\begin{advdefinition}{牛顿流体}\label{advA1:def:newtonian-fluid}
采用“压应力为负”的张量符号时，
\[
 \bm\sigma=-p\mathsf1+2\eta\bm U'
 +\zeta(\nabla\cdot\vect v)\mathsf1,
\]
其中 $\eta$ 是动力黏度，$\zeta$ 是体黏度（第二黏度）。
\end{advdefinition}

\begin{advdefinition}{正压流体}\label{advA1:def:barotropic}
若物态方程可写为 $\rho=\rho(p)$，或在指定熵、温度条件下把密度视为压强的单值函数，则称为正压流体。此时
$\nabla\rho\times\nabla p=0$。
\end{advdefinition}

\begin{advdefinition}{理想流体}\label{advA1:def:ideal-fluid}
忽略全部黏滞偏应力，即
\[
 \bm\sigma=-p\mathsf1.
\]
\end{advdefinition}

\begin{advdefinition}{不可压缩流体}\label{advA1:def:incompressible}
沿质元 $\rho$ 为常数；对均匀密度情形，
\[
 \rho=\text{常量},\qquad \nabla\cdot\vect v=0.
\]
\end{advdefinition}

\begin{advtheorem}{理想流体的欧拉方程}\label{advA1:thm:euler-fluid}
无体力时，动量守恒的守恒形式和随体形式分别是
\[
 \frac{\partial(\rho\vect v)}{\partial t}
 +\nabla\cdot(\rho\vect v\vect v+p\mathsf1)=0,
\]
\[
 \boxed{\rho\frac{D\vect v}{Dt}=-\nabla p}.
\]
\end{advtheorem}
证明：对任意固定控制体应用动量守恒。对流动量通量是
$\rho\vect v\vect v$，表面压强力用散度定理化为 $-\nabla p$。再用质量连续性方程消去
$\vect v[\partial_t\rho+\nabla\cdot(\rho\vect v)]$，便得到随体形式。二次通量使欧拉方程本质上是非线性的。

\begin{advtheorem}{开尔文环流定理}\label{advA1:thm:kelvin}
对无黏、正压且体力为保守力的流动，任意随体闭曲线 $C(t)$ 上的环流
\[
 \Gamma=\oint_{C(t)}\vect v\cdot\dd\vect l
\]
满足 $D\Gamma/Dt=0$。
\end{advtheorem}
证明：随体线元满足
$D(\dd\vect l)/Dt=(\dd\vect l\cdot\nabla)\vect v$，故
\[
 \frac{D\Gamma}{Dt}=\oint\left[
 \frac{D\vect v}{Dt}\cdot\dd\vect l+
 \vect v\cdot(\dd\vect l\cdot\nabla)\vect v\right].
\]
第二项是 $\oint\dd(v^2/2)=0$。正压流体中
$\nabla p/\rho=\nabla h(p)$，保守体力也是梯度，闭路积分均为零。

\begin{advexamplebox}{涡度方程}\label{advA1:ex:vorticity}
在无黏正压流体中证明
\[
 \frac{D\vect\omega}{Dt}
 =(\vect\omega\cdot\nabla)\vect v
 -(\nabla\cdot\vect v)\vect\omega .
\]
\end{advexamplebox}
\begin{advsolution}
对欧拉方程取旋度。正压条件使
$\nabla\times(\nabla p/\rho)=
\nabla(1/\rho)\times\nabla p=0$。利用
\[
 (\vect v\cdot\nabla)\vect v
 =\nabla\frac{v^2}{2}-\vect v\times\vect\omega
\]
得
\[
 \partial_t\vect\omega=\nabla\times(\vect v\times\vect\omega).
\]
展开叉积旋度，并用 $\nabla\cdot\vect\omega=0$：
\[
 \partial_t\vect\omega
 +(\vect v\cdot\nabla)\vect\omega
 =(\vect\omega\cdot\nabla)\vect v
 -(\nabla\cdot\vect v)\vect\omega.
\]
左边即随体导数。不可压缩流体中第二项消失；二维无伸展流动还会使涡量随质元保持不变。
\end{advsolution}

\begin{advexamplebox}{应变率张量的随体演化}\label{advA1:ex:strain-rate}
设理想流体满足 $\rho=\rho(p)$。证明
\[
 \frac{D\bm U}{Dt}
 =-\frac{\nabla\nabla p}{\rho}
 +\frac{\dd\rho}{\dd p}\frac{\nabla p\otimes\nabla p}{\rho^2}
 -\bm U^2-\bm W^2.
\]
\end{advexamplebox}
\begin{advsolution}
记 $\mathsf A=\nabla\vect v=\bm U+\bm W$。对速度梯度取随体导数时，随体导数与空间梯度不对易：
\[
 \frac{D\mathsf A}{Dt}
 =\nabla\left(\frac{D\vect v}{Dt}\right)-\mathsf A^2.
\]
欧拉方程给出
\[
 -\nabla\left(\frac{\nabla p}{\rho}\right)
 =-\frac{\nabla\nabla p}{\rho}
 +\frac{\dd\rho}{\dd p}\frac{\nabla p\otimes\nabla p}{\rho^2}.
\]
取矩阵的对称部分。因
$\operatorname{sym}(\mathsf A^2)=\bm U^2+\bm W^2$，交叉项
$\bm U\bm W+\bm W\bm U$ 是反对称的，故得到所证公式。
\end{advsolution}

若初始无旋且满足开尔文定理的条件，涡度方程没有产生涡量的源项，流动保持无旋，可写
$\vect v=\nabla\Phi$。利用
\[
 (\vect v\cdot\nabla)\vect v
 =\nabla(v^2/2)-\vect v\times(\nabla\times\vect v),
\]
正压流体的欧拉方程可积分为
\[
 \partial_t\Phi+\frac{v^2}{2}+\int^p\frac{\dd p'}{\rho(p')}=C(t).
\]

\begin{advexamplebox}{半圆柱形机库的升力}\label{advA1:ex:hangar-lift}
半圆柱形机库截面是上半圆 $r=R$，沿 $z$ 方向无限长。外部不可压缩无旋空气的速度势为
\[
 \Phi=-v_0\left(r+\frac{R^2}{r}\right)\cos\theta.
\]
若机库内部压强为远处静压 $p_\infty$，求单位长度机库受到的竖直升力。
\end{advexamplebox}
\begin{advsolution}
速度分量为
\[
 v_r=\partial_r\Phi=-v_0\left(1-\frac{R^2}{r^2}\right)\cos\theta,
\]
\[
 v_\theta=\frac1r\partial_\theta\Phi
 =v_0\left(1+\frac{R^2}{r^2}\right)\sin\theta.
\]
在表面 $r=R$，$v_r=0$、$v_\theta=2v_0\sin\theta$。伯努利方程给出外表面压强
\[
 p_{\rm out}=p_\infty+\frac12\rho
 (v_0^2-4v_0^2\sin^2\theta).
\]
内外压强差沿外法线产生的竖直分力为
\[
 \frac{F_y}{H}
 =R\int_0^\pi(p_\infty-p_{\rm out})\sin\theta\,\dd\theta
 =\frac12\rho v_0^2R\left(
 4\int_0^\pi\sin^3\theta\,\dd\theta-\int_0^\pi\sin\theta\,\dd\theta\right).
\]
用 $\int_0^\pi\sin^3\theta\,\dd\theta=4/3$，得
\[
 \boxed{\frac{F_y}{H}=\frac53\rho v_0^2R}.
\]
\begin{center}\begin{minipage}{\linewidth}\centering
\begin{tikzpicture}[scale=.92]
 \draw[ground] (-2.2,0)--(2.2,0);
 \draw[very thick] (-1.65,0) arc[start angle=180,end angle=0,radius=1.65];
 \foreach \a in {20,50,90,130,160}{
  \draw[vector]
    ({1.65*cos(\a)-.28*sin(\a)},{1.65*sin(\a)+.28*cos(\a)})
    --({1.65*cos(\a)+.28*sin(\a)},{1.65*sin(\a)-.28*cos(\a)});
 }
 \draw[vector] (-2.5,1.45)--(-1.6,1.45) node[midway,above]{$v_0$};
 \draw[axis] (0,0)--(0,2.15) node[above]{$y$};
 \node[fill=white] at (0,.55) {机库内部 $p_\infty$};
\end{tikzpicture}
\bookcaption{半圆柱表面速度增大导致外压降低，压力差的竖直分量形成升力}\label{advA1:fig:hangar}
\end{minipage}\end{center}
\end{advsolution}

热力学基本关系对随体质元写成
\[
 \frac{Du_s}{Dt}=-p\frac{Dv_s}{Dt}+T\frac{Ds_s}{Dt},
 \qquad v_s=1/\rho.
\]
准静态绝热、无耗散时 $Ds_s/Dt=0$，内能变化与速度散度耦合。

\begin{advtheorem}{理想流体的能量方程}\label{advA1:thm:bernoulli-energy}
令总能量密度
\[
 w=\rho u_s+\frac12\rho v^2.
\]
无体力、无黏且准静态绝热时，
\[
 \boxed{\frac{\partial w}{\partial t}
 +\nabla\cdot[(w+p)\vect v]=0}.
\]
\end{advtheorem}
证明：用欧拉方程点乘 $\vect v$ 得到宏观动能方程；用
$Du_s/Dt=-(p/\rho)\nabla\cdot\vect v$ 与质量连续性方程得到内能方程。两式相加，压强功整理为
$\nabla\cdot(p\vect v)$，即得守恒形式。

\begin{advexamplebox}{伯努利原理的两种形式}\label{advA1:ex:bernoulli-principle}
(1) 理想流体受单位质量势能 $\phi(\vect r)$，证明准静态绝热定常流中沿流线
\[
 u_s+\frac p\rho+\frac{v^2}{2}+\phi=C.
\]
(2) 对不可压缩理想流体，证明即使不使用准静态绝热条件，也有
$p/\rho+v^2/2+\phi$ 沿流线不变；若流动无旋，则常数在整个连通流区相同。
\end{advexamplebox}
\begin{advsolution}
(1) 把定常能量方程写成随体形式，并用热力学恒等式
$\dd h=T\dd s+\dd p/\rho$，其中 $h=u_s+p/\rho$。沿绝热流线
$\dd s=0$，故
\[
 \vect v\cdot\nabla\left(
 h+\frac{v^2}{2}+\phi\right)=0.
\]

(2) 不可压缩欧拉方程是
\[
 (\vect v\cdot\nabla)\vect v
 =-\nabla\left(\frac p\rho+\phi\right).
\]
与 $\vect v$ 点乘：
\[
 \vect v\cdot\nabla\left(
 \frac{v^2}{2}+\frac p\rho+\phi\right)=0,
\]
所以沿流线守恒。若 $\nabla\times\vect v=0$，则
$(\vect v\cdot\nabla)\vect v=\nabla(v^2/2)$，整个括号的梯度为零，故在连通流区全局为同一常数。
\end{advsolution}

\begin{advtheorem}{达朗贝尔佯谬}\label{advA1:thm:dalembert}
在无环流、不可压缩、无黏、定常势流模型中，匀速来流对封闭物体的合阻力为零。
\end{advtheorem}
证明可用包围物体的远处控制面，不要求物体前后对称。
对三维局域封闭物体，假定流体在无穷远趋于均匀来流且无源汇、无尾迹。
势流的远场扰动从偶极项开始，故
$\vect v=\vect U+\vect u$ 中 $\vect u=O(R^{-3})$。
Bernoulli 关系给出 $p-p_\infty=-\rho\vect U\cdot\vect u+O(u^2)$。
稳恒动量通量用 $p\mathsf1+\rho\vect v\vect v$ 表示；
均匀部分在闭合控制面上积分为零，扰动部分至多为
$O(R^2R^{-3})\to0$。物面不可穿透，因而总压强力也为零。
二维无环流的局域势流可作相应远场论证；若有环流，仍需区分升力与阻力。
前后几何对称只允许对特殊形状配对压强，不能作为一般证明。
现实阻力来自黏性边界层、分离和尾迹，恰是此理想模型删除的机制。

\subsubsection{不可压缩牛顿流体}

简单平行剪切中
\[
 \tau=\eta\frac{\dd v}{\dd y}.
\]
动力黏度单位为 $\si{Pa.s}$，运动黏度 $\nu=\eta/\rho$ 的单位为
$\si{m^2.s^{-1}}$。
\begin{center}\begin{minipage}{\linewidth}\centering
\begin{tikzpicture}[scale=.9]
 \draw[ground] (-.2,0)--(4.8,0);
 \draw[ground] (-.2,2.2)--(4.8,2.2);
 \foreach \y/\len in {.4/.7,.8/1.25,1.2/1.8,1.6/2.35,2.0/2.9}{
  \draw[vector] (.45,\y)--++(\len,0);
 }
 \draw[<->,TextGray] (4.2,.8)--(4.2,1.2) node[midway,right]{$\dd y$};
 \draw[guide] (1.7,.8)--(4.3,.8) (2.25,1.2)--(4.3,1.2);
\node at (2.2,2.52) {上板带动流体};
\end{tikzpicture}
\bookcaption{牛顿黏滞定律：相邻流层的速度差产生剪应力}\label{advA1:fig:newton-viscosity}
\end{minipage}\end{center}
图\ref{advA1:fig:newton-viscosity}中箭头长度表示速度，标签放在流层外侧，没有遮住速度线。

不可压缩牛顿流体满足纳维--斯托克斯方程
\[
 \rho\left(\frac{\partial\vect v}{\partial t}
 +\vect v\cdot\nabla\vect v\right)
 =-\nabla p+\eta\nabla^2\vect v+\vect f.
\]
可压缩情形还会出现体黏度与 $\nabla(\nabla\cdot\vect v)$ 项。圆管泊肃叶定律
$\Delta p=8\eta QL/(\pi R^4)$ 和低雷诺数球体斯托克斯阻力
$F=6\pi\eta Rv$ 都是其经典特解或渐近结果。

\begin{advexamplebox}{同轴圆柱间的旋转 Couette 流}\label{advA1:ex:cylindrical-couette}
高为 $H$ 的同轴圆柱半径为 $R_1<R_2$，角速度分别为
$\omega_1,\omega_2$。间隙内充满不可压缩牛顿流体，黏度为 $\eta$。求定常速度场及维持两圆柱转动所需外力矩。
\end{advexamplebox}
\begin{advsolution}
轴对称定常解只有方位分量，纳维--斯托克斯方程化为
\[
 \frac{\dd^2v_\theta}{\dd r^2}
 +\frac1r\frac{\dd v_\theta}{\dd r}
 -\frac{v_\theta}{r^2}=0,
\]
故
\[
 v_\theta=Ar+\frac Br,
\]
\[
 A=\frac{\omega_2R_2^2-\omega_1R_1^2}{R_2^2-R_1^2},
 \qquad
 B=\frac{R_1^2R_2^2(\omega_1-\omega_2)}
 {R_2^2-R_1^2}.
\]
剪应力为
\[
 \tau_{r\theta}=\eta\left(
 \frac{\dd v_\theta}{\dd r}-\frac{v_\theta}{r}\right)
 =-\frac{2\eta B}{r^2}.
\]
任一圆柱面上传递的力矩大小与半径无关：
\[
 |\mathcal M|=2\pi rH\cdot r|\tau_{r\theta}|
 =\boxed{4\pi\eta H|B|}.
\]
若 $\omega_1>\omega_2$，维持内筒需沿正方向施加
$\mathcal M_1=4\pi\eta HB$，维持外筒需施加
$\mathcal M_2=-4\pi\eta HB$；两者之和为零，功率差转化为黏滞热。
\begin{center}\begin{minipage}{\linewidth}\centering
\begin{tikzpicture}[scale=.88]
 \draw[very thick] (0,0) circle(.85);
 \draw[very thick] (0,0) circle(1.65);
 \draw[->,DeepBlue,thick] (.85,0) arc[start angle=0,end angle=105,radius=.85];
 \draw[->,TextGray,thick] (1.65,0) arc[start angle=0,end angle=55,radius=1.65];
 \draw[<->,TextGray] (0,0)--(.85,0) node[midway,below]{$R_1$};
 \draw[<->,TextGray] (0,0)--({1.65*cos(-35)},{1.65*sin(-35)}) node[midway,below]{$R_2$};
 \node at (.30,.56) {$\omega_1$};
 \node at (1.45,1.25) {$\omega_2$};
\end{tikzpicture}
\bookcaption{同轴圆柱之间的方位 Couette 流与两侧无滑移边界}\label{advA1:fig:couette}
\end{minipage}\end{center}
\end{advsolution}

\begin{advexamplebox}{环隙中稳定下落圆柱测黏度}\label{advA1:ex:falling-cylinder-viscosity}
半径 $R$ 的圆柱杯内盛密度 $\rho_0$、黏度未知的液体。半径 $r<R$、高 $a$、密度
$\rho>\rho_0$ 的小圆柱与杯同轴，以稳定速度 $v$ 下落。圆柱完全浸没，杯底封闭，液体不可压缩；取远离杯底的长环隙、低雷诺数极限，忽略端部局部损失，但保留连续性所要求的压差与回流。求黏度。
\end{advexamplebox}
\begin{advsolution}
取向下为 $z$ 正向，设圆柱完全浸没，液体不可压缩且杯底封闭。
记 $H=\ln(R/r)$、$D=R^2-r^2$，把静水压梯度扣除后令
$G=\dd p/\dd z-\rho_0g$。充分发展方程与无滑移边界为
\[
 \eta\frac1s\frac{\dd}{\dd s}\left(s\frac{\dd u}{\dd s}\right)=G,
 \qquad u(r)=v,\quad u(R)=0.
\]
积分得到剪切驱动与压差驱动的叠加：
\[
 u(s)=v\frac{\ln(R/s)}H+
 \frac G{4\eta}\left[s^2-R^2+\frac D H\ln\frac R s\right].
\]
圆柱每秒排开 $\pi r^2v$ 的液体，故环隙必须向上回流：
\[
 2\pi\int_r^R u(s)s\,\dd s=-\pi r^2v.
\]
代入速度分布，令 $K=H(R^2+r^2)-D>0$，得
\[
 G=\frac{4\eta v}K.
\]
除侧面剪力外，端面还受到相对于静水压的压差，二者都是向上阻力：
\[
 F_{\rm drag}=-2\pi ra\eta u'(r)+G\pi r^2a
 =\frac{2\pi\eta av(R^2+r^2)}K.
\]
把静水浮力另计，终端速度条件为
$(\rho-\rho_0)\pi r^2ag=F_{\rm drag}$，从而
\[
 \boxed{\eta=\frac{(\rho-\rho_0)r^2g}{2v}
 \left[\ln\frac Rr-\frac{R^2-r^2}{R^2+r^2}\right]}.
\]
方括号严格为正：令 $R/r=\e^H$，它就是 $H-\tanh H>0$。
窄隙 $h=R-r\ll r$ 时约为 $h^3/(3r^3)$，于是
$\eta\simeq(\rho-\rho_0)gh^3/(6rv)$。
原来的纯对数速度剖面仅有剪切流，会给出向下净流量，不满足封底杯的连续性条件。
本式保留压差但忽略入口、出口的局部损失，适用于长圆柱、低雷诺数和远离杯底的充分发展段；
“忽略端部绕流细节”不等于可以删去端面压差。
\begin{center}\begin{minipage}{\linewidth}\centering
\begin{tikzpicture}[scale=.88]
 \draw[very thick] (-2,-2)--(-2,2);
 \draw[very thick] (2,-2)--(2,2);
 \draw[ground] (-2,-2)--(2,-2);
 \fill[DeepBlue!6] (-2,-1.95) rectangle (2,1.15);
 \draw[very thick,fill=white] (-.75,-.65) rectangle (.75,1.05);
 \draw[vector] (0,.25)--(0,-.60) node[midway,right]{$v$};
 \draw[<->,TextGray] (-.75,1.38)--(.75,1.38) node[midway,above]{$2r$};
 \draw[<->,TextGray] (-2,-1.55)--(2,-1.55) node[midway,below]{$2R$};
 \draw[<->,TextGray] (1.05,-.65)--(1.05,1.05) node[midway,right]{$a$};
\end{tikzpicture}
\bookcaption{同轴下落圆柱与杯壁之间的轴向环隙层流}\label{advA1:fig:falling-cylinder}
\end{minipage}\end{center}
图\ref{advA1:fig:falling-cylinder}明确画出 $r<R$、共轴和侧面环隙；液体在环隙内向上净回流，端部局部损失按长环隙近似忽略。
\end{advsolution}

}

\begin{advexamplebox}{等周不等式}\label{advA1:ex:isoperimetric}
设光滑闭曲线长度为 $C$，证明其围成的有向面积满足
\[
 |A|=\left|\frac12\oint(x\,\dd y-y\,\dd x)\right|
 \le\frac{C^2}{4\pi},
\]
并指出等号条件。
\end{advexamplebox}
\begin{advsolution}
以弧长 $s\in[0,C]$ 参数化，并平移原点使
$\int_0^C x\,\dd s=\int_0^C y\,\dd s=0$。周期函数的维廷格不等式给出
\[
 \int_0^C(x^2+y^2)\dd s\le
 \left(\frac{C}{2\pi}\right)^2
 \int_0^C(x'^2+y'^2)\dd s.
\]
弧长参数满足 $x'^2+y'^2=1$。再用柯西不等式：
\[
 2|A|=\left|\int_0^C(xy'-yx')\dd s\right|
 \le\sqrt{\int_0^C(x^2+y^2)\dd s}
 \sqrt{\int_0^C(x'^2+y'^2)\dd s}
 \le\frac{C^2}{2\pi}.
\]
因此 $|A|\le C^2/(4\pi)$。两次不等式同时取等号要求
$x,y$ 只含第一傅里叶谐波且切矢量处处与位矢垂直，故曲线必须是圆。
\end{advsolution}

\begin{advexamplebox}{最速降线}\label{advA1:ex:brachistochrone}
重力沿 $y$ 轴正向，质点从 $(0,0)$ 静止释放，沿光滑轨道到达 $(l,h)$。忽略摩擦，求使下降时间
\[
 \mathcal T[y]=\int_0^l\sqrt{\frac{1+y'^2}{2gy}}\,\dd x
\]
最小的曲线。
\end{advexamplebox}
\begin{advsolution}
密度函数不显含 $x$，可用贝尔特拉米恒等式
\[
 \mathcal L-y'\frac{\partial\mathcal L}{\partial y'}
 =\frac1{\sqrt{2gy(1+y'^2)}}=C_0 .
\]
令 $y=a(1-\cos\theta)$，则由上式得到
\[
 x=a(\theta-\sin\theta)+C_1.
\]
取起点为 $\theta=0$：
\[
 \boxed{x=a(\theta-\sin\theta),\qquad
 y=a(1-\cos\theta)}.
\]
这是半径为 $a$ 的圆滚动时圆周上一点描出的旋轮线。参数 $a$ 与终点参数 $\theta_1$ 由
$l=a(\theta_1-\sin\theta_1)$、$h=a(1-\cos\theta_1)$ 唯一确定适用分支。
\begin{center}\begin{minipage}{\linewidth}\centering
\begin{tikzpicture}[x=1.05cm,y=.85cm]
 \draw[axis] (0,0)--(4.2,0) node[right]{$x$};
 \draw[axis] (0,0)--(0,-2.65) node[below]{$y$（向下）};
 \draw[curve,domain=0.03:2.8,samples=120,smooth,variable=\t]
   plot ({\t-sin(\t r)},{-(1-cos(\t r))});
 \fill (0,0) circle(2pt) node[above left]{$O$};
 \fill ({2.8-sin(2.8 r)},{-(1-cos(2.8 r))}) circle(2pt) node[right]{$P$};
 \draw[guide] (0,{-(1-cos(2.8 r))})--({2.8-sin(2.8 r)},{-(1-cos(2.8 r))});
\end{tikzpicture}
\bookcaption{最速降线是旋轮线，而不是直线或圆弧}\label{advA1:fig:brachistochrone}
\end{minipage}\end{center}
图\ref{advA1:fig:brachistochrone}中轨道在起点具有竖直切线，使质点先迅速获得速度，再用较高速度完成横向运动。
\end{advsolution}

\begin{advexamplebox}{最小耗散原理}\label{advA1:ex:least-dissipation}
区域 $D$ 内电导率为 $\sigma(\vect r)>0$，给定非静电力场 $\vect K$。在固定非静电力功率
\[
 -\int_D \vect K\cdot\sigma\nabla\phi\,\dd V=P
\]
的条件下，求使给定二次泛函
$Q[\phi]=\int_D\sigma|\nabla\phi|^2\dd V$ 取极值的函数，并解释它与稳恒导电方程的形式关系，辨明何时可解释为物理耗散。
\end{advexamplebox}
\begin{advsolution}
对
\[
 \mathcal F[\phi]=\int_D\sigma|\nabla\phi|^2\dd V
 +\lambda\left(\int_D\vect K\cdot\sigma\nabla\phi\,\dd V+P\right)
\]
变分。若边界电势固定，则
\[
 \delta\mathcal F
 =-\int_D\nabla\cdot\{\sigma(2\nabla\phi+\lambda\vect K)\}
 \,\delta\phi\,\dd V .
\]
所以
\[
 \nabla\cdot\{\sigma(2\nabla\phi+\lambda\vect K)\}=0.
\]
把电势作常数重标度 $V=2\phi/(-\lambda)$，便得到
\[
 \nabla\cdot\vect j=0,\qquad
 \vect j=\sigma(\vect K-\nabla V),
\]
即欧姆定律与稳恒电荷守恒的形式；重标度也改变边界电势，且只在 $\lambda\ne0$ 时有效。二阶变分
$\delta^2Q=2\int_D\sigma|\nabla\delta\phi|^2\dd V>0$，所以在固定边界、线性约束的可行集合内这是最小值。$\lambda$ 仍须由给定的 $P$ 求定；一般没有 $Q_{\min}=P$。例如零边值的一维区间 $[0,1]$，取无量纲 $\sigma=1,K=x$，则约束为 $\int_0^1\phi\,\dd x=P$，极小解为 $\phi=6Px(1-x)$，且 $Q_{\min}=12P^2$。

对于重标度后的物理电流，真正的焦耳功率是
\[
 Q_J=\int_D\frac{j^2}{\sigma}\dd V
 =\int_D\sigma|\vect K-\nabla V|^2\dd V,
\]
并非原泛函 $Q[\phi]$。由连续性可得
\[
 Q_J=\int_D\vect K\cdot\vect j\,\dd V
 -\oint_{\partial D}V\vect j\cdot\dd\vect S.
\]
只有边界电功率项为零时，非静电力的输入功率才等于焦耳热。
\end{advsolution}

\begin{advexamplebox}{中心力场中的最速降线}\label{advA1:ex:central-brachistochrone}
质点从半径 $r_0$ 处静止释放。分别对
(1) 平方反比引力场；(2) 均匀球形天体内部的线性引力场，写出最速降线的确定方程。
\end{advexamplebox}
\begin{advsolution}
若速率只依赖 $r$，下降时间是光学型泛函
\[
 \mathcal T=\int n(r)\,\dd s,\qquad n(r)=\frac1{v(r)}.
\]
旋转对称性给出守恒量
\[
 b=n(r)r\sin\alpha,
\]
其中 $\alpha$ 是轨道切线与径向的夹角。因此所有中心力最速降线均由
\begin{equation}\label{advA1:eq:central-brach-quad}
 \frac{\dd\varphi}{\dd r}
 =\frac{b}{r\sqrt{n^2(r)r^2-b^2}}
\end{equation}
以及端点条件确定。参数 $b$ 的物理意义与分层介质中的斯涅耳常量相同。

(1) 对 $U=-GMm/r$，
\[
 v(r)=\sqrt{2GM\left(\frac1r-\frac1{r_0}\right)},\qquad
 n(r)=\left[2GM\left(\frac1r-\frac1{r_0}\right)\right]^{-1/2}.
\]
代入式\eqref{advA1:eq:central-brach-quad}得到径向单调段；对终点 $r_1<r_0$ 且无径向转折的候选轨道，用角位移积分
\[
 \Delta\varphi=\int_{r_1}^{r_0}
 \frac{b\,\dd r}{r\sqrt{r^2/[2GM(1/r-1/r_0)]-b^2}}
\]
求 $b$，再由 $\mathcal T=\int n\,\dd s$ 求时间。

(2) 均匀球内 $g(r)=\Omega_g^2r$，$\Omega_g^2=GM/R^3$。从 $r_0$ 静止释放时
\[
 v(r)=\Omega_g\sqrt{r_0^2-r^2},\qquad
 n(r)=\frac1{\Omega_g\sqrt{r_0^2-r^2}}.
\]
令 $q=\Omega_g b$，轨道满足
\[
 \frac{\dd\varphi}{\dd r}
 =\frac{q\sqrt{r_0^2-r^2}}
 {r\sqrt{(1+q^2)r^2-q^2r_0^2}}.
\]
上述导数的正负号须随径向运动方向选择。根号要求轨道不越过由
$n(r)r=b$ 决定的最近心点 $r_*$。若先向内后向外，角位移须改为
\[
 |\Delta\varphi|=\int_{r_*}^{r_0}F(r;b)\,\dd r
 +\int_{r_*}^{r_1}F(r;b)\,\dd r,
 \quad F(r;b)=\frac{|b|}{r\sqrt{n^2(r)r^2-b^2}}.
\]
同半径异角端点取 $r_1=r_0$ 时，两段贡献相加而不相消。角条件可能给出多个驻时轨道，还须比较正的总时间并检查轨道是否位于所用引力模型的区域；不能把欧拉方程的一个积分分支直接宣称为唯一最速轨道。
\end{advsolution}

对力学轨道 $q_\alpha(t)$，作用量为
\[
 S[q]=\int_{t_1}^{t_2}L(q,\dot q,t)\dd t.
\]

\begin{advtheorem}{哈密顿原理}\label{advA1:thm:hamilton-principle}
在固定起止时刻与端点的条件下，真实轨道使作用量的一阶变分为零：
\[
 \delta S=0.
\]
它给出的欧拉--拉格朗日方程正是分析力学的拉格朗日方程。
\end{advtheorem}
“一阶变分为零”比“全局最小”更准确，因为真实轨道也可能对应局部最大或鞍点。

\begin{advexamplebox}{曲面上的测地线}\label{advA1:ex:geodesic}
曲面 $\vect r=\vect r(u^1,u^2)$ 的线元为
\[
 \dd s^2=g_{ij}(u)\dd u^i\dd u^j.
\]
求参数曲线 $u^i=u^i(t)$ 为测地线的条件。
\end{advexamplebox}
\begin{advsolution}
若 $t$ 是仿射参数，可对能量泛函
\[
 S=\frac12\int g_{ij}\dot u^i\dot u^j\dd t
\]
变分。欧拉--拉格朗日方程为
\[
 \frac{\dd}{\dd t}(g_{kj}\dot u^j)
 -\frac12\partial_k g_{ij}\dot u^i\dot u^j=0.
\]
左乘逆度规 $g^{\ell k}$，得
\[
 \boxed{\ddot u^\ell+\Gamma^\ell{}_{ij}\dot u^i\dot u^j=0},
\qquad
 \Gamma^\ell{}_{ij}=\frac12g^{\ell k}
 (\partial_i g_{kj}+\partial_j g_{ki}-\partial_k g_{ij}).
\]
若用任意参数直接变分弧长，则右边可出现
$f(t)\dot u^\ell$；重新参数化为弧长或其他仿射参数后该项消失。
\end{advsolution}

\begin{advexamplebox}{极小曲面的条件}\label{advA1:ex:minimal-surface}
曲面 $\vect r=\vect r(u^1,u^2)$ 是某空间区域的边界。求固定边界曲线时它使面积取驻值的条件，并写出图形曲面 $z=f(x,y)$ 的方程。
\end{advexamplebox}
\begin{advsolution}
面积泛函为
\[
 A[\vect r]=\int\sqrt{\det g}\,\dd u^1\dd u^2.
\]
沿法向作变分 $\delta\vect r=\eta\vect n$，一阶面积变分为
\[
 \delta A=-2\int H\eta\,\dd A+\oint_{\partial S}(\cdots),
\]
固定边界使线积分消失。由于 $\eta$ 任意，极小曲面的必要条件为平均曲率
\[
 \boxed{H=0}.
\]
对 $z=f(x,y)$，
\[
 A[f]=\iint\sqrt{1+f_x^2+f_y^2}\,\dd x\dd y.
\]
代入欧拉--拉格朗日方程，得到
\[
 \boxed{(1+f_y^2)f_{xx}-2f_xf_yf_{xy}
 +(1+f_x^2)f_{yy}=0}.
\]
平面显然满足该方程；悬链面等非平凡曲面也满足，但球面因 $H\ne0$ 不是无约束极小曲面。
\end{advsolution}

\subsection{弹性体与场论守恒律}

先看一维弹性棒。位移场为 $\xi(x,t)$，质量线密度与等效弹性系数记为 $\rho,E$：
\[
 \mathcal L=\frac12\rho(\partial_t\xi)^2-\frac12E(\partial_x\xi)^2.
\]
场方程、波速与能量连续性方程分别为
\[
 \rho\,\partial_t^2\xi-E\,\partial_x^2\xi=0,\qquad
 c=\sqrt{E/\rho},
\]
\[
 w=\frac12\rho\xi_t^2+\frac12E\xi_x^2,\qquad
 \Pi=-E\xi_t\xi_x,\qquad
 \partial_t w+\partial_x\Pi=0.
\]
相应的机械动量密度与动量流为
\[
 g=\rho\xi_t,\qquad \mathcal T=-E\xi_x,\qquad
 \partial_tg+\partial_x\mathcal T=0.
\]

\begin{advtheorem}{场论的诺特定理}\label{advA1:thm:field-noether}
若标量场 $\phi_\nu$ 的连续变换
$\bar\phi_\nu=\Lambda_\nu(\phi,\lambda)$ 在 $\lambda=0$ 的生成元为
$\Psi_\nu=(\partial\Lambda_\nu/\partial\lambda)_0$，且拉格朗日密度的变化是散度
\[
 \left.\frac{\partial\bar{\mathcal L}}{\partial\lambda}\right|_0
 =\partial_\mu B^\mu,
\]
则
\[
 j^\mu=B^\mu-\sum_\nu
 \frac{\partial\mathcal L}{\partial(\partial_\mu\phi_\nu)}\Psi_\nu,
 \qquad \partial_\mu j^\mu=0.
\]
\end{advtheorem}
证明：把 $\partial_\lambda\mathcal L$ 用链式法则展开，对含
$\partial_\mu\Psi_\nu$ 的项作乘积求导，再用场方程
$\partial_\mu[\partial\mathcal L/\partial(\partial_\mu\phi_\nu)]
=\partial\mathcal L/\partial\phi_\nu$，余项正好整理为 $\partial_\mu j^\mu=0$。

若 $\mathcal L$ 不显含时空坐标，四个时空平移参数给出能量--动量张量
\[
 T^\mu{}_\nu=
 \sum_a\frac{\partial\mathcal L}{\partial(\partial_\mu\phi_a)}
 \partial_\nu\phi_a-\delta^\mu{}_\nu\mathcal L,
 \qquad \partial_\mu T^\mu{}_\nu=0.
\]
其时间分量是能量密度与能流，空间分量是动量密度与动量流。

\begin{advexamplebox}{弹性棒的平移对称性与机械动量}\label{advA1:ex:rod-noether}
证明一维弹性棒的拉格朗日密度不显含 $\xi$ 所对应的连续对称性产生机械动量守恒。
\end{advexamplebox}
\begin{advsolution}
变换 $\xi\mapsto\xi+\epsilon$ 不改变 $\xi_t,\xi_x$，故 $\delta\mathcal L=0$，生成元为 $\Psi=1$。
诺特流可整体乘以 $-1$ 重新约定为
\[
 j^0=\frac{\partial\mathcal L}{\partial\xi_t}=\rho\xi_t,
 \qquad
 j^1=\frac{\partial\mathcal L}{\partial\xi_x}=-E\xi_x.
\]
连续性方程
\[
 \partial_t(\rho\xi_t)+\partial_x(-E\xi_x)=0
\]
正是波动方程，也就是机械动量密度 $g=\rho\xi_t$ 与动量流
$\mathcal T=-E\xi_x$ 的守恒式。
\end{advsolution}

三维弹性体用拉格朗日描述
$\vect R(\vect r,t)=\vect r+\vect\xi(\vect r,t)$。变形后的基矢
$\vect\varepsilon_i=\partial\vect R/\partial x_i$ 构成度规
$g_{ij}=\vect\varepsilon_i\cdot\vect\varepsilon_j$。

\begin{advdefinition}{应变张量}\label{advA1:def:strain}
\[
 \varepsilon_{ij}=\frac12(g_{ij}-\delta_{ij}).
\]
在线性小形变下，对角元是正应变，非对角元是工程剪应变的一半。
\end{advdefinition}

\begin{advexamplebox}{纯剪切把单位圆变成椭圆}\label{advA1:ex:shear-circle}
二维小应变张量为
\[
 \bm\varepsilon=\begin{pmatrix}0&\theta/2\\\theta/2&0\end{pmatrix}.
\]
求单位圆 $x^2+y^2=1$ 的变形后曲线，并对角化应变张量。
\end{advexamplebox}
\begin{advsolution}
取不含刚性转动的代表变形梯度
$\mathsf F=\mathsf1+\bm\varepsilon$：
\[
 x'=x+\frac\theta2y,\qquad
 y'=y+\frac\theta2x.
\]
保留到 $\theta$ 一阶，逆变换为
$x=x'-\theta y'/2,\ y=y'-\theta x'/2$，所以
\[
 \boxed{x'^2+y'^2-2\theta x'y'=1+\order{\theta^2}}.
\]
取旋转坐标
$u=(x+y)/\sqrt2,\ v=(x-y)/\sqrt2$，有
\[
 \mathsf O\bm\varepsilon\mathsf O^{\mathsf T}
 =\operatorname{Diag}(\theta/2,-\theta/2).
\]
故精确到线性阶，主轴方向的伸长率分别为 $1+\theta/2$ 与
$1-\theta/2$：
\[
 \frac{u'^2}{(1+\theta/2)^2}
 +\frac{v'^2}{(1-\theta/2)^2}=1.
\]
\begin{center}\begin{minipage}{\linewidth}\centering
\begin{tikzpicture}[scale=1.0]
 \draw[axis] (-2.0,0)--(2.0,0) node[right]{$x'$};
 \draw[axis] (0,-1.55)--(0,1.55) node[above]{$y'$};
 \draw[guide] (0,0) circle(1.0);
 \begin{scope}[rotate=45]
  \draw[curve] (0,0) ellipse (1.36 and .74);
  \draw[guide] (-1.55,0)--(1.55,0) node[right]{$u$};
  \draw[guide] (0,-1.05)--(0,1.05) node[above]{$v$};
 \end{scope}
\end{tikzpicture}
\bookcaption{纯剪应变的主轴沿 $x=\pm y$，单位圆变为转过 $45^\circ$ 的椭圆}\label{advA1:fig:shear-ellipse}
\end{minipage}\end{center}
图\ref{advA1:fig:shear-ellipse}中虚线圆仅作变形前参照，标签没有用白底遮挡曲线。
\end{advsolution}

\begin{advtheorem}{应变张量的位移场表示}\label{advA1:thm:strain-displacement}
在线性小形变近似下，
\[
 \boxed{\bm\varepsilon=\frac12[
 \nabla\vect\xi+(\nabla\vect\xi)^{\mathsf T}]}.
\]
\end{advtheorem}
证明：$\mathsf F=\mathsf1+\nabla\vect\xi$，故
\[
 \mathsf g=\mathsf F^{\mathsf T}\mathsf F
 =\mathsf1+\nabla\vect\xi+(\nabla\vect\xi)^{\mathsf T}
 +(\nabla\vect\xi)^{\mathsf T}\nabla\vect\xi.
\]
忽略位移梯度的二阶项，再用
$\bm\varepsilon=(\mathsf g-\mathsf1)/2$ 即得。反对称部分
\[
 \bm\varphi=\frac12[
 \nabla\vect\xi-(\nabla\vect\xi)^{\mathsf T}]
\]
只描述质元的刚性小转动，不改变线元长度。

\begin{advtheorem}{$SO(3)$ 张量积的不可约分解}\label{advA1:thm:so3-decomposition}
三维旋转群的整数角动量不可约实表示维数为 $2\ell+1$。对
$m=2\ell_1+1\ge n=2\ell_2+1$，
\[
 m\otimes n=(m+n-1)\oplus(m+n-3)\oplus\cdots\oplus(m-n+1).
\]
特别地，
\[
 3\otimes3=5\oplus3\oplus1.
\]
\end{advtheorem}
对二阶张量 $\mathsf P$，这三个部分分别是
\[
 \mathsf P_{\rm ST}
 =\frac{\mathsf P+\mathsf P^{\mathsf T}}2
 -\frac{\operatorname{tr}\mathsf P}{3}\mathsf1,\quad
 \mathsf P_{\rm A}=\frac{\mathsf P-\mathsf P^{\mathsf T}}2,\quad
 \mathsf P_{\rm tr}=\frac{\operatorname{tr}\mathsf P}{3}\mathsf1.
\]
它们的维数依次为 $5,3,1$。对称无迹应变描述定体积形变，迹
$\operatorname{tr}\bm\varepsilon$ 是体积相对变化，反对称部分是刚性转动。

\begin{advexamplebox}{三阶张量的不可约分解计数}\label{advA1:ex:rank3-decomp}
证明维数分解
\[
 3\otimes3\otimes3=7\oplus2\cdot5\oplus3\cdot3\oplus1.
\]
\end{advexamplebox}
\begin{advsolution}
先耦合前两个矢量表示：
\[
 3\otimes3=5\oplus3\oplus1
 \quad(\ell=2,1,0).
\]
再分别与第三个 $\ell=1$ 的矢量耦合：
\[
 5\otimes3=7\oplus5\oplus3,\qquad
 3\otimes3=5\oplus3\oplus1,\qquad
 1\otimes3=3.
\]
合并同类项即得
$7\oplus2\cdot5\oplus3\cdot3\oplus1$。维数检查
$7+10+9+1=27=3^3$，没有遗漏分量。
\end{advsolution}

\AOneElasticSecond
\AOneFluidSection


\clearpage\endgroup
